\documentclass[../main.tex]{subfiles}
\begin{document}

\chapter{Remote Procedure Calls}
\lessoninfo{
  video={P4L1},
  textbook={OSTEP \cite{ostep} ch.~48; Silberschatz et al.\ \cite{silberschatz2018} ch.~3; Birrell and Nelson \cite{birrell1984}; RFC 5531 \cite{rfc5531}, RFC 4506 \cite{rfc4506}},
  keywords={remote procedure call, stub, RPC runtime, interface definition language, marshalling, unmarshalling, binding, registry, partial failure, Sun RPC, XDR, rpcgen, rpcbind, Java RMI},
}

The preceding lessons considered a single machine, managed by one
operating system or, in Lesson~12, by a hypervisor beneath several.  This
lesson opens the part on distributed systems, in which cooperating
processes run on different machines connected by a network.  Across
machines, the processes can only exchange messages, typically through the
sockets of Lesson~9, and every application that uses them repeats the same
low-level communication code.  Remote procedure calls hide this code behind
the interface of an ordinary procedure call.  The lesson presents the
structure of RPC systems and the design decisions that each of them must
make---interface definition, marshalling, binding, pointer arguments and
partial failures---and then examines how Sun RPC and Java RMI make them.

\section{Why remote procedure calls?}
\label{sec:l13-why}

Consider two client-server applications.\source{P4L1, why RPC, benefits of
RPC; Birrell and Nelson \cite{birrell1984}; OSTEP ch.~48.}  In the first, a
client requests a file from a server, which returns the contents of the
file.  In the second, a client sends an image together with the name and
the parameters of an image-processing algorithm, and the server returns the
processed image.  The services differ, yet the communication code of the
two applications is largely the same.  Each of them must
\begin{itemize}
  \item create and initialize the sockets through which client and server
    communicate;
  \item allocate buffers for the messages and populate them;
  \item include protocol information that tells the other side how to
    interpret a message, such as the operation requested (retrieve a file,
    apply an algorithm) and the size of the data that follow;
  \item copy the data into the buffers: a file name and the contents of the
    file in the first application, the algorithm parameters and the image
    in the second;
\end{itemize}
and the receiving side performs the inverse steps.  Only the operations and
the types of the data exchanged are specific to the application.

\begin{definition}[Remote procedure call]
  A \term{remote procedure call}[遠隔手続き呼び出し] (RPC) is an
  inter-process communication mechanism with which a process invokes a
  procedure that executes in another address space, usually on another
  machine, in the same way as it invokes a local procedure.  The RPC system
  performs the communication that transfers the arguments to the procedure
  and the results back to the caller.
\end{definition}
\index{RPC|see{remote procedure call}}

RPC is intended to simplify the development of interactions that cross
address spaces and machines.  Its benefits are the following.
\begin{description}
  \item[Higher-level interface] Data movement and communication are
    expressed as calls of procedures with typed arguments and results,
    rather than as byte buffers written to and read from sockets.
  \item[Error handling] The RPC system detects communication errors and
    reports them to the application in a uniform way.
  \item[Hiding complexity] The details of the interaction between machines,
    from connection management to differences in how the machines represent
    data, are handled by the RPC system and by code that it generates.
\end{description}

Birrell and Nelson \cite{birrell1984} described the design and
implementation of an RPC facility for the Cedar programming environment at
Xerox PARC in 1984.  The structure they chose, with generated stubs and a
runtime library on both sides (\cref{sec:l13-structure}), underlies later
RPC systems, including the two examined in this lesson.

\section{RPC requirements}
\label{sec:l13-requirements}

The goal of making a remote operation look like a procedure call leads to
the following requirements for an RPC system.\source{P4L1, RPC
requirements; OSTEP ch.~48.}
\begin{description}
  \item[Client-server interactions] RPC targets the client-server model: a
    server offers a set of procedures, and clients call them.
  \item[Procedure call interface] An RPC has \emph{synchronous call
    semantics}.  The calling process, or in a multithreaded process the
    calling thread, blocks when it makes the call and continues only when
    the procedure has completed and its result has been returned, exactly
    as for a local call.
  \item[Type checking] Because arguments and results have declared types,
    the RPC system can check them and report an error when a call does not
    match the procedure.  The types also tell the receiver how to interpret
    the bytes of a message, which a socket delivers as an uninterpreted
    sequence.
  \item[Cross-machine conversion] Client and server may run on machines that
    represent data differently, and the RPC system must convert between
    the representations.
  \item[Higher-level protocol] RPC is a protocol above the transport layer.
    It can incorporate access control, authenticating clients and
    authorizing their calls, and mechanisms for fault tolerance, and it
    should be able to use different transport protocols, such as TCP and
    UDP.
\end{description}

The most common difference in representation is the order in which the
bytes of a multi-byte value are stored.  A machine is
\term{big-endian}[ビッグエンディアン] if it stores the most significant byte
of an integer at the lowest address, and
\term{little-endian}[リトルエンディアン] if it stores the least significant
byte there.  The x86 processors are little-endian; many network protocols,
and the encoding of \cref{sec:l13-encoding}, use big-endian order.

\begin{example}
  \label{ex:l13-endian}
  The 32-bit integer $300 = \mathtt{0x0000012C}$ occupies the bytes
  \texttt{00 00 01 2C}, in order of increasing address, on a big-endian
  machine and \texttt{2C 01 00 00} on a little-endian machine.  If a
  little-endian client copied the bytes of this integer into a message
  unchanged, a big-endian server would read them as
  $\mathtt{0x2C010000} = \num{738263040}$.
\end{example}

\section{Structure of an RPC system}
\label{sec:l13-structure}

In a local procedure call, the caller places the arguments on the stack or
in registers and jumps to the procedure, which computes the result and
returns.\source{P4L1, structure of RPC, RPC steps; Birrell and Nelson
\cite{birrell1984}; OSTEP ch.~48.}  In an RPC, caller and procedure reside in
different address spaces, so neither the arguments nor the code of the
procedure are accessible to the other side.  The call therefore jumps to a
local procedure that has the same interface as the remote one but performs
the communication instead (\cref{fig:l13-steps}).

\begin{definition}[Stub]
  A \term{stub}[スタブ] is a procedure, generated by the RPC system, that
  converts between a procedure call and messages.  The \emph{client stub},
  linked into the client, has the interface of the remote procedure: it
  packs the arguments into a message and unpacks the result from the reply.
  The \emph{server stub}, linked into the server, unpacks the arguments,
  calls the actual implementation of the procedure and packs its result
  into the reply.
\end{definition}

\begin{definition}[RPC runtime]
  The \term{RPC runtime}[RPCランタイム] is the library, present on both
  sides, that transmits the messages between client and server: it manages
  the connections, sends and receives the messages, retransmits lost ones,
  and passes each received message to the right stub.
\end{definition}

Neither caller nor procedure contains any communication code; each sees
an ordinary local call.\margin{Birrell and Nelson's system had the same
parts, with a stub generator called Lupine \cite{birrell1984}.}

\begin{figure}[htbp]
  \centering
  % Structure of an RPC system and the steps of one call.
% Badges: -1 register, 0 bind, 1 call, 2 marshal, 3 send, 4 receive,
% 5 unmarshal, 6 actual call, 7 result.  Dashed arrows: return path.
% Usage: % Structure of an RPC system and the steps of one call.
% Badges: -1 register, 0 bind, 1 call, 2 marshal, 3 send, 4 receive,
% 5 unmarshal, 6 actual call, 7 result.  Dashed arrows: return path.
% Usage: % Structure of an RPC system and the steps of one call.
% Badges: -1 register, 0 bind, 1 call, 2 marshal, 3 send, 4 receive,
% 5 unmarshal, 6 actual call, 7 result.  Dashed arrows: return path.
% Usage: \input{figures/l13-steps}   (width ~116 mm)
\begin{tikzpicture}[x=1mm, y=1mm, font=\footnotesize\sffamily,
  >={Stealth[length=1.6mm]},
  machine/.style={draw=ln-rule, rounded corners=2pt, fill=ln-box},
  comp/.style={draw, fill=white, align=center, minimum width=34mm,
               minimum height=8mm, inner sep=1pt},
  stubc/.style={comp, fill=ln-accent!15},
  rt/.style={comp, fill=ln-accent!30},
  reg/.style={draw, rounded corners=2pt, fill=white, align=center,
              minimum width=26mm, minimum height=7mm, inner sep=1pt},
  stepn/.style={circle, draw, fill=white, inner sep=0pt,
                minimum size=4.2mm, font=\scriptsize\sffamily\bfseries},
  hdr/.style={font=\scriptsize\sffamily\bfseries, text=ln-muted},
  lab/.style={font=\scriptsize\sffamily, text=ln-muted},
  ret/.style={->, densely dashed}]
  % machines
  \draw[machine] (2,0) rectangle (46,48);
  \draw[machine] (74,0) rectangle (118,48);
  \node[hdr] at (24,2.6) {\iflangja{クライアントマシン}{client machine}};
  \node[hdr] at (96,2.6) {\iflangja{サーバマシン}{server machine}};
  % client side
  \node[comp]  (caller) at (24,38)
    {\iflangja{呼び出し元}{caller}\\[-0.3ex]\texttt{s = sum(3, v);}};
  \node[stubc] (cstub)  at (24,24.5) {\iflangja{クライアントスタブ}{client stub}};
  \node[rt]    (crt)    at (24,11) {\iflangja{RPC ランタイム}{RPC runtime}};
  % server side
  \node[comp]  (proc)   at (96,38)
    {\iflangja{手続きの実装}{implementation}\\[-0.3ex]\texttt{sum(count, values)}};
  \node[stubc] (sstub)  at (96,24.5) {\iflangja{サーバスタブ}{server stub}};
  \node[rt]    (srt)    at (96,11) {\iflangja{RPC ランタイム}{RPC runtime}};
  % registry
  \node[reg] (reg) at (60,56) {\iflangja{レジストリ}{registry}};
  % -1 register, 0 bind
  \draw[->] ([xshift=-6mm]proc.north) |- (reg.east);
  \draw[<->] ([xshift=6mm]caller.north) |- (reg.west);
  \node[stepn] at (90,51.5) {$-1$};
  \node[stepn] at (30,51.5) {0};
  % client: call path down, return path up
  \draw[->]  ([xshift=-8mm]caller.south) -- ([xshift=-8mm]cstub.north);
  \draw[->]  ([xshift=-8mm]cstub.south) -- ([xshift=-8mm]crt.north);
  \draw[ret] ([xshift=8mm]crt.north) -- ([xshift=8mm]cstub.south);
  \draw[ret] ([xshift=8mm]cstub.north) -- ([xshift=8mm]caller.south);
  \node[stepn] at (11.5,31.3) {1};
  \node[stepn] at (7,24.5) {2};
  % network
  \draw[->]  ([yshift=2mm]crt.east) -- ([yshift=2mm]srt.west);
  \draw[ret] ([yshift=-2mm]srt.west) -- ([yshift=-2mm]crt.east);
  \node[stepn] at (51,17) {3};
  \node[stepn] at (69,17) {4};
  \node[lab] at (60,5.5) {\iflangja{ネットワーク}{network}};
  % server: request path up, return path down
  \draw[->]  ([xshift=-8mm]srt.north) -- ([xshift=-8mm]sstub.south);
  \draw[->]  ([xshift=-8mm]sstub.north) -- ([xshift=-8mm]proc.south);
  \draw[ret] ([xshift=8mm]proc.south) -- ([xshift=8mm]sstub.north);
  \draw[ret] ([xshift=8mm]sstub.south) -- ([xshift=8mm]srt.north);
  \node[stepn] at (113,24.5) {5};
  \node[stepn] at (83.5,31.3) {6};
  \node[stepn] at (113,38) {7};
\end{tikzpicture}
   (width ~116 mm)
\begin{tikzpicture}[x=1mm, y=1mm, font=\footnotesize\sffamily,
  >={Stealth[length=1.6mm]},
  machine/.style={draw=ln-rule, rounded corners=2pt, fill=ln-box},
  comp/.style={draw, fill=white, align=center, minimum width=34mm,
               minimum height=8mm, inner sep=1pt},
  stubc/.style={comp, fill=ln-accent!15},
  rt/.style={comp, fill=ln-accent!30},
  reg/.style={draw, rounded corners=2pt, fill=white, align=center,
              minimum width=26mm, minimum height=7mm, inner sep=1pt},
  stepn/.style={circle, draw, fill=white, inner sep=0pt,
                minimum size=4.2mm, font=\scriptsize\sffamily\bfseries},
  hdr/.style={font=\scriptsize\sffamily\bfseries, text=ln-muted},
  lab/.style={font=\scriptsize\sffamily, text=ln-muted},
  ret/.style={->, densely dashed}]
  % machines
  \draw[machine] (2,0) rectangle (46,48);
  \draw[machine] (74,0) rectangle (118,48);
  \node[hdr] at (24,2.6) {\iflangja{クライアントマシン}{client machine}};
  \node[hdr] at (96,2.6) {\iflangja{サーバマシン}{server machine}};
  % client side
  \node[comp]  (caller) at (24,38)
    {\iflangja{呼び出し元}{caller}\\[-0.3ex]\texttt{s = sum(3, v);}};
  \node[stubc] (cstub)  at (24,24.5) {\iflangja{クライアントスタブ}{client stub}};
  \node[rt]    (crt)    at (24,11) {\iflangja{RPC ランタイム}{RPC runtime}};
  % server side
  \node[comp]  (proc)   at (96,38)
    {\iflangja{手続きの実装}{implementation}\\[-0.3ex]\texttt{sum(count, values)}};
  \node[stubc] (sstub)  at (96,24.5) {\iflangja{サーバスタブ}{server stub}};
  \node[rt]    (srt)    at (96,11) {\iflangja{RPC ランタイム}{RPC runtime}};
  % registry
  \node[reg] (reg) at (60,56) {\iflangja{レジストリ}{registry}};
  % -1 register, 0 bind
  \draw[->] ([xshift=-6mm]proc.north) |- (reg.east);
  \draw[<->] ([xshift=6mm]caller.north) |- (reg.west);
  \node[stepn] at (90,51.5) {$-1$};
  \node[stepn] at (30,51.5) {0};
  % client: call path down, return path up
  \draw[->]  ([xshift=-8mm]caller.south) -- ([xshift=-8mm]cstub.north);
  \draw[->]  ([xshift=-8mm]cstub.south) -- ([xshift=-8mm]crt.north);
  \draw[ret] ([xshift=8mm]crt.north) -- ([xshift=8mm]cstub.south);
  \draw[ret] ([xshift=8mm]cstub.north) -- ([xshift=8mm]caller.south);
  \node[stepn] at (11.5,31.3) {1};
  \node[stepn] at (7,24.5) {2};
  % network
  \draw[->]  ([yshift=2mm]crt.east) -- ([yshift=2mm]srt.west);
  \draw[ret] ([yshift=-2mm]srt.west) -- ([yshift=-2mm]crt.east);
  \node[stepn] at (51,17) {3};
  \node[stepn] at (69,17) {4};
  \node[lab] at (60,5.5) {\iflangja{ネットワーク}{network}};
  % server: request path up, return path down
  \draw[->]  ([xshift=-8mm]srt.north) -- ([xshift=-8mm]sstub.south);
  \draw[->]  ([xshift=-8mm]sstub.north) -- ([xshift=-8mm]proc.south);
  \draw[ret] ([xshift=8mm]proc.south) -- ([xshift=8mm]sstub.north);
  \draw[ret] ([xshift=8mm]sstub.south) -- ([xshift=8mm]srt.north);
  \node[stepn] at (113,24.5) {5};
  \node[stepn] at (83.5,31.3) {6};
  \node[stepn] at (113,38) {7};
\end{tikzpicture}
   (width ~116 mm)
\begin{tikzpicture}[x=1mm, y=1mm, font=\footnotesize\sffamily,
  >={Stealth[length=1.6mm]},
  machine/.style={draw=ln-rule, rounded corners=2pt, fill=ln-box},
  comp/.style={draw, fill=white, align=center, minimum width=34mm,
               minimum height=8mm, inner sep=1pt},
  stubc/.style={comp, fill=ln-accent!15},
  rt/.style={comp, fill=ln-accent!30},
  reg/.style={draw, rounded corners=2pt, fill=white, align=center,
              minimum width=26mm, minimum height=7mm, inner sep=1pt},
  stepn/.style={circle, draw, fill=white, inner sep=0pt,
                minimum size=4.2mm, font=\scriptsize\sffamily\bfseries},
  hdr/.style={font=\scriptsize\sffamily\bfseries, text=ln-muted},
  lab/.style={font=\scriptsize\sffamily, text=ln-muted},
  ret/.style={->, densely dashed}]
  % machines
  \draw[machine] (2,0) rectangle (46,48);
  \draw[machine] (74,0) rectangle (118,48);
  \node[hdr] at (24,2.6) {\iflangja{クライアントマシン}{client machine}};
  \node[hdr] at (96,2.6) {\iflangja{サーバマシン}{server machine}};
  % client side
  \node[comp]  (caller) at (24,38)
    {\iflangja{呼び出し元}{caller}\\[-0.3ex]\texttt{s = sum(3, v);}};
  \node[stubc] (cstub)  at (24,24.5) {\iflangja{クライアントスタブ}{client stub}};
  \node[rt]    (crt)    at (24,11) {\iflangja{RPC ランタイム}{RPC runtime}};
  % server side
  \node[comp]  (proc)   at (96,38)
    {\iflangja{手続きの実装}{implementation}\\[-0.3ex]\texttt{sum(count, values)}};
  \node[stubc] (sstub)  at (96,24.5) {\iflangja{サーバスタブ}{server stub}};
  \node[rt]    (srt)    at (96,11) {\iflangja{RPC ランタイム}{RPC runtime}};
  % registry
  \node[reg] (reg) at (60,56) {\iflangja{レジストリ}{registry}};
  % -1 register, 0 bind
  \draw[->] ([xshift=-6mm]proc.north) |- (reg.east);
  \draw[<->] ([xshift=6mm]caller.north) |- (reg.west);
  \node[stepn] at (90,51.5) {$-1$};
  \node[stepn] at (30,51.5) {0};
  % client: call path down, return path up
  \draw[->]  ([xshift=-8mm]caller.south) -- ([xshift=-8mm]cstub.north);
  \draw[->]  ([xshift=-8mm]cstub.south) -- ([xshift=-8mm]crt.north);
  \draw[ret] ([xshift=8mm]crt.north) -- ([xshift=8mm]cstub.south);
  \draw[ret] ([xshift=8mm]cstub.north) -- ([xshift=8mm]caller.south);
  \node[stepn] at (11.5,31.3) {1};
  \node[stepn] at (7,24.5) {2};
  % network
  \draw[->]  ([yshift=2mm]crt.east) -- ([yshift=2mm]srt.west);
  \draw[ret] ([yshift=-2mm]srt.west) -- ([yshift=-2mm]crt.east);
  \node[stepn] at (51,17) {3};
  \node[stepn] at (69,17) {4};
  \node[lab] at (60,5.5) {\iflangja{ネットワーク}{network}};
  % server: request path up, return path down
  \draw[->]  ([xshift=-8mm]srt.north) -- ([xshift=-8mm]sstub.south);
  \draw[->]  ([xshift=-8mm]sstub.north) -- ([xshift=-8mm]proc.south);
  \draw[ret] ([xshift=8mm]proc.south) -- ([xshift=8mm]sstub.north);
  \draw[ret] ([xshift=8mm]sstub.south) -- ([xshift=8mm]srt.north);
  \node[stepn] at (113,24.5) {5};
  \node[stepn] at (83.5,31.3) {6};
  \node[stepn] at (113,38) {7};
\end{tikzpicture}

  \caption{Structure of an RPC system and the steps of a call.  Solid
    arrows: the call; dashed arrows: the return of the result.  Steps $-1$
    and~0 take place before the first call.}
  \label{fig:l13-steps}
\end{figure}

\subsection{Steps of an RPC}

An RPC proceeds in the following steps, numbered as in
\cref{fig:l13-steps}.
\begin{description}
  \item[$-1$ Register] The server registers the procedures it offers, their
    argument types and its location, so that clients can find it
    (\cref{sec:l13-binding}).
  \item[0 Bind] The client finds a server that provides the desired
    procedures and binds to it, for example by opening a connection.
  \item[1 Call] The client calls the procedure.  Control passes to the
    client stub, and the caller blocks.
  \item[2 Marshal] The client stub marshals the arguments: it serializes
    them into a buffer (\cref{sec:l13-marshalling}).
  \item[3 Send] The RPC runtime of the client sends the message to the
    server, with the transport protocol agreed on during binding.
  \item[4 Receive] The RPC runtime of the server receives the message and
    passes it to the server stub.  Access control checks are performed at
    this point.
  \item[5 Unmarshal] The server stub unmarshals the arguments: it extracts
    them from the buffer and creates the corresponding data structures.
  \item[6 Actual call] The server stub calls the local implementation of the
    procedure.
  \item[7 Result] The procedure computes the result, which returns to the
    client along the reverse path: the server stub marshals it, the
    runtimes transfer the reply, the client stub unmarshals it and returns
    it to the caller, which resumes.
\end{description}
Steps $-1$ and~0 are performed once, before the first call; steps 1--7 are
repeated for every call.

\begin{keypoint}
  An RPC looks like a local procedure call to both the caller and the
  procedure.  The stubs convert between calls and messages, and the RPC
  runtime moves the messages between the machines.
\end{keypoint}

\section{Interface definition language}
\label{sec:l13-idl}

Client and server must agree on what the server can do and on the
arguments and results of each of its procedures.\source{P4L1, interface
definition language, specifying an IDL; OSTEP ch.~48.}  The client needs
this information to decide whether to bind to a server, and the RPC system
needs it in a precise, machine-readable form to generate the stubs
automatically.

\begin{definition}[Interface definition language]
  An \term{interface definition language}[インタフェース定義言語] (IDL) is a
  language for describing the interface that a server exports: the names
  of its procedures, the types of their arguments and results, and a
  version number.
\end{definition}
\index{IDL|see{interface definition language}}

An RPC system can use one of two kinds of IDL.
\begin{description}
  \item[Language-agnostic IDL] A language of its own, independent of the
    programming languages in which clients and servers are written, from
    which stubs are generated for each supported language.  Sun RPC uses
    XDR (\cref{sec:l13-sunrpc}).
  \item[Language-specific IDL] The interface is written in the programming
    language of the application itself, so programmers need not learn
    another language.  Java RMI uses Java (\cref{sec:l13-rmi}).
\end{description}
In either case the IDL describes only the interface of the server, not its
implementation.

The version number supports the evolution of a service.  When a service
is upgraded, the servers are often not all upgraded at the same time, and a
server may export the old and the new version of an interface side by
side.  By stating the version it expects, a client binds only to servers
that implement that version.

\section{Marshalling and unmarshalling}
\label{sec:l13-marshalling}

The arguments of a call reside in the memory of the caller: values on the
stack and in registers, and data elsewhere in the address space to which
pointers refer.\source{P4L1, marshalling, unmarshalling; OSTEP ch.~48.}
A message, in contrast, is a contiguous sequence of bytes.

\begin{definition}[Marshalling]
  The conversion of the arguments or the result of a call into a
  contiguous buffer, in an encoding agreed between client and server, so
  that they can be sent as a message and correctly interpreted by the
  receiver, is called \term{marshalling}[マーシャリング] or
  \emph{serialization}.
\end{definition}
\index{serialization|see{marshalling}}

\begin{definition}[Unmarshalling]
  The inverse conversion is called \term{unmarshalling}[アンマーシャリング]:
  parsing a received buffer according to the types given by the interface,
  extracting the values, and creating the corresponding data structures in
  the address space of the receiver, allocating and initializing memory as
  needed.
\end{definition}

For data of fixed size, the types in the interface tell the receiver where
one value ends and the next begins.  Data of variable size, such as arrays
and strings, need additional information: either their length is encoded
before their elements, or their end is marked by a special value, as the
NUL character terminates a C string.

\begin{example}
  \label{ex:l13-marshal}
  The interface of a server declares a procedure
  \texttt{int sum(int count, int *values)} and states that
  \texttt{values} points to \texttt{count} integers.  The client calls
  \texttt{sum(3, v)}, where \texttt{v} is an array at address
  \texttt{0x5a10} that holds 7, $-2$ and 300.  With
  \SI{4}{\byte} per integer, the client stub builds the \SI{16}{\byte}
  buffer of \cref{fig:l13-marshal}: the count followed by the three
  elements.  The address \texttt{0x5a10} is not sent.  The server stub reads
  the count, allocates memory for three integers, here at \texttt{0x9c40},
  copies the elements into it, and calls the implementation with
  \texttt{count} $=3$ and \texttt{values} $=\mathtt{0x9c40}$.
\end{example}

\begin{figure}[htbp]
  \centering
  % Marshalling and unmarshalling the arguments of sum(count, values).
% The pointer value is not sent; the pointed-to array is copied into the
% buffer and recreated at a different address on the server.
% Usage: % Marshalling and unmarshalling the arguments of sum(count, values).
% The pointer value is not sent; the pointed-to array is copied into the
% buffer and recreated at a different address on the server.
% Usage: % Marshalling and unmarshalling the arguments of sum(count, values).
% The pointer value is not sent; the pointed-to array is copied into the
% buffer and recreated at a different address on the server.
% Usage: \input{figures/l13-marshal}   (width ~117 mm, fits the 120 mm body)
\begin{tikzpicture}[x=1mm, y=1mm, font=\footnotesize\sffamily,
  >={Stealth[length=1.6mm]},
  cell/.style={draw, fill=white, minimum height=5mm, inner sep=0pt,
               font=\scriptsize\sffamily},
  var/.style={cell, minimum width=16mm},
  elem/.style={cell, minimum width=7mm},
  buf/.style={cell, minimum width=7mm, fill=ln-accent!15},
  vname/.style={font=\scriptsize\sffamily, anchor=east},
  hdr/.style={font=\scriptsize\sffamily\bfseries, text=ln-muted, anchor=base},
  lab/.style={font=\scriptsize\sffamily, text=ln-muted},
  addr/.style={font=\scriptsize\ttfamily, text=ln-muted, anchor=south west,
               inner sep=0.5pt},
  op/.style={font=\scriptsize\sffamily, anchor=south}]
  % ---------------- client address space
  \node[hdr] at (15,42) {\iflangja{クライアントのアドレス空間}{client address space}};
  \node[lab, anchor=west] at (0,36.5) {\iflangja{スタック}{stack}};
  \node[var] (cc) at (21,31) {3};
  \node[var] (cv) at (21,25) {\texttt{0x5a10}};
  \node[vname] at (12.5,31) {\texttt{count}};
  \node[vname] at (12.5,25) {\texttt{values}};
  \node[elem] (c1) at (11.5,9) {7};
  \node[elem] (c2) at (18.5,9) {$-2$};
  \node[elem] (c3) at (25.5,9) {300};
  \node[addr, anchor=north west] at (8,6.3) {0x5a10};
  \node[lab, anchor=east] at (7.6,9) {\iflangja{ヒープ}{heap}};
  \draw[->] ([xshift=-4mm]cv.south) -- (c1.north);
  % ---------------- message buffer
  \node[hdr] at (60,42) {\iflangja{メッセージバッファ}{message buffer}};
  \foreach \v [count=\i from 0] in {3,7,$-2$,300} {
    \node[buf] (b\i) at (49.5+7*\i,20) {\v};
    \node[font=\tiny\sffamily, text=ln-muted] at (46+7*\i,24.5) {\the\numexpr4*\i\relax};
  }
  \node[font=\tiny\sffamily, text=ln-muted] at (74,24.5) {16};
  \draw[ln-muted] (46.4,16.8) -- (46.4,15.8) -- (52.6,15.8) -- (52.6,16.8);
  \node[lab, anchor=north] at (49.5,15.3) {\texttt{count}};
  \draw[ln-muted] (53.4,16.8) -- (53.4,15.8) -- (73.6,15.8) -- (73.6,16.8);
  \node[lab, anchor=north] at (63.5,15.3)
    {\iflangja{要素}{elements}};
  % ---------------- server address space
  \node[hdr] at (103,42) {\iflangja{サーバのアドレス空間}{server address space}};
  \node[lab, anchor=west] at (88,36.5) {\iflangja{スタック}{stack}};
  \node[var] (sc) at (109,31) {3};
  \node[var] (sv) at (109,25) {\texttt{0x9c40}};
  \node[vname] at (100.5,31) {\texttt{count}};
  \node[vname] at (100.5,25) {\texttt{values}};
  \node[elem] (s1) at (99.5,9) {7};
  \node[elem] (s2) at (106.5,9) {$-2$};
  \node[elem] (s3) at (113.5,9) {300};
  \node[addr, anchor=north west] at (96,6.3) {0x9c40};
  \node[lab, anchor=east] at (95.6,9) {\iflangja{ヒープ}{heap}};
  \draw[->] ([xshift=-4mm]sv.south) -- (s1.north);
  % ---------------- operations
  % Japanese: below the arrows; above them the wider "アンマーシャル"
  % runs into the buffer offset "16".
  \draw[->, thick, ln-accent] (31,20) -- (44.5,20);
  \node[op, anchor=\iflangja{north}{south}] at (37.8,\iflangja{19.2}{20.8}) {\iflangja{マーシャル}{marshal}};
  \draw[->, thick, ln-accent] (75.5,20) -- (87,20);
  \node[op, anchor=\iflangja{north}{south}] at (\iflangja{82.3}{81.3},\iflangja{19.2}{20.8}) {\iflangja{アンマーシャル}{unmarshal}};
\end{tikzpicture}
   (width ~117 mm, fits the 120 mm body)
\begin{tikzpicture}[x=1mm, y=1mm, font=\footnotesize\sffamily,
  >={Stealth[length=1.6mm]},
  cell/.style={draw, fill=white, minimum height=5mm, inner sep=0pt,
               font=\scriptsize\sffamily},
  var/.style={cell, minimum width=16mm},
  elem/.style={cell, minimum width=7mm},
  buf/.style={cell, minimum width=7mm, fill=ln-accent!15},
  vname/.style={font=\scriptsize\sffamily, anchor=east},
  hdr/.style={font=\scriptsize\sffamily\bfseries, text=ln-muted, anchor=base},
  lab/.style={font=\scriptsize\sffamily, text=ln-muted},
  addr/.style={font=\scriptsize\ttfamily, text=ln-muted, anchor=south west,
               inner sep=0.5pt},
  op/.style={font=\scriptsize\sffamily, anchor=south}]
  % ---------------- client address space
  \node[hdr] at (15,42) {\iflangja{クライアントのアドレス空間}{client address space}};
  \node[lab, anchor=west] at (0,36.5) {\iflangja{スタック}{stack}};
  \node[var] (cc) at (21,31) {3};
  \node[var] (cv) at (21,25) {\texttt{0x5a10}};
  \node[vname] at (12.5,31) {\texttt{count}};
  \node[vname] at (12.5,25) {\texttt{values}};
  \node[elem] (c1) at (11.5,9) {7};
  \node[elem] (c2) at (18.5,9) {$-2$};
  \node[elem] (c3) at (25.5,9) {300};
  \node[addr, anchor=north west] at (8,6.3) {0x5a10};
  \node[lab, anchor=east] at (7.6,9) {\iflangja{ヒープ}{heap}};
  \draw[->] ([xshift=-4mm]cv.south) -- (c1.north);
  % ---------------- message buffer
  \node[hdr] at (60,42) {\iflangja{メッセージバッファ}{message buffer}};
  \foreach \v [count=\i from 0] in {3,7,$-2$,300} {
    \node[buf] (b\i) at (49.5+7*\i,20) {\v};
    \node[font=\tiny\sffamily, text=ln-muted] at (46+7*\i,24.5) {\the\numexpr4*\i\relax};
  }
  \node[font=\tiny\sffamily, text=ln-muted] at (74,24.5) {16};
  \draw[ln-muted] (46.4,16.8) -- (46.4,15.8) -- (52.6,15.8) -- (52.6,16.8);
  \node[lab, anchor=north] at (49.5,15.3) {\texttt{count}};
  \draw[ln-muted] (53.4,16.8) -- (53.4,15.8) -- (73.6,15.8) -- (73.6,16.8);
  \node[lab, anchor=north] at (63.5,15.3)
    {\iflangja{要素}{elements}};
  % ---------------- server address space
  \node[hdr] at (103,42) {\iflangja{サーバのアドレス空間}{server address space}};
  \node[lab, anchor=west] at (88,36.5) {\iflangja{スタック}{stack}};
  \node[var] (sc) at (109,31) {3};
  \node[var] (sv) at (109,25) {\texttt{0x9c40}};
  \node[vname] at (100.5,31) {\texttt{count}};
  \node[vname] at (100.5,25) {\texttt{values}};
  \node[elem] (s1) at (99.5,9) {7};
  \node[elem] (s2) at (106.5,9) {$-2$};
  \node[elem] (s3) at (113.5,9) {300};
  \node[addr, anchor=north west] at (96,6.3) {0x9c40};
  \node[lab, anchor=east] at (95.6,9) {\iflangja{ヒープ}{heap}};
  \draw[->] ([xshift=-4mm]sv.south) -- (s1.north);
  % ---------------- operations
  % Japanese: below the arrows; above them the wider "アンマーシャル"
  % runs into the buffer offset "16".
  \draw[->, thick, ln-accent] (31,20) -- (44.5,20);
  \node[op, anchor=\iflangja{north}{south}] at (37.8,\iflangja{19.2}{20.8}) {\iflangja{マーシャル}{marshal}};
  \draw[->, thick, ln-accent] (75.5,20) -- (87,20);
  \node[op, anchor=\iflangja{north}{south}] at (\iflangja{82.3}{81.3},\iflangja{19.2}{20.8}) {\iflangja{アンマーシャル}{unmarshal}};
\end{tikzpicture}
   (width ~117 mm, fits the 120 mm body)
\begin{tikzpicture}[x=1mm, y=1mm, font=\footnotesize\sffamily,
  >={Stealth[length=1.6mm]},
  cell/.style={draw, fill=white, minimum height=5mm, inner sep=0pt,
               font=\scriptsize\sffamily},
  var/.style={cell, minimum width=16mm},
  elem/.style={cell, minimum width=7mm},
  buf/.style={cell, minimum width=7mm, fill=ln-accent!15},
  vname/.style={font=\scriptsize\sffamily, anchor=east},
  hdr/.style={font=\scriptsize\sffamily\bfseries, text=ln-muted, anchor=base},
  lab/.style={font=\scriptsize\sffamily, text=ln-muted},
  addr/.style={font=\scriptsize\ttfamily, text=ln-muted, anchor=south west,
               inner sep=0.5pt},
  op/.style={font=\scriptsize\sffamily, anchor=south}]
  % ---------------- client address space
  \node[hdr] at (15,42) {\iflangja{クライアントのアドレス空間}{client address space}};
  \node[lab, anchor=west] at (0,36.5) {\iflangja{スタック}{stack}};
  \node[var] (cc) at (21,31) {3};
  \node[var] (cv) at (21,25) {\texttt{0x5a10}};
  \node[vname] at (12.5,31) {\texttt{count}};
  \node[vname] at (12.5,25) {\texttt{values}};
  \node[elem] (c1) at (11.5,9) {7};
  \node[elem] (c2) at (18.5,9) {$-2$};
  \node[elem] (c3) at (25.5,9) {300};
  \node[addr, anchor=north west] at (8,6.3) {0x5a10};
  \node[lab, anchor=east] at (7.6,9) {\iflangja{ヒープ}{heap}};
  \draw[->] ([xshift=-4mm]cv.south) -- (c1.north);
  % ---------------- message buffer
  \node[hdr] at (60,42) {\iflangja{メッセージバッファ}{message buffer}};
  \foreach \v [count=\i from 0] in {3,7,$-2$,300} {
    \node[buf] (b\i) at (49.5+7*\i,20) {\v};
    \node[font=\tiny\sffamily, text=ln-muted] at (46+7*\i,24.5) {\the\numexpr4*\i\relax};
  }
  \node[font=\tiny\sffamily, text=ln-muted] at (74,24.5) {16};
  \draw[ln-muted] (46.4,16.8) -- (46.4,15.8) -- (52.6,15.8) -- (52.6,16.8);
  \node[lab, anchor=north] at (49.5,15.3) {\texttt{count}};
  \draw[ln-muted] (53.4,16.8) -- (53.4,15.8) -- (73.6,15.8) -- (73.6,16.8);
  \node[lab, anchor=north] at (63.5,15.3)
    {\iflangja{要素}{elements}};
  % ---------------- server address space
  \node[hdr] at (103,42) {\iflangja{サーバのアドレス空間}{server address space}};
  \node[lab, anchor=west] at (88,36.5) {\iflangja{スタック}{stack}};
  \node[var] (sc) at (109,31) {3};
  \node[var] (sv) at (109,25) {\texttt{0x9c40}};
  \node[vname] at (100.5,31) {\texttt{count}};
  \node[vname] at (100.5,25) {\texttt{values}};
  \node[elem] (s1) at (99.5,9) {7};
  \node[elem] (s2) at (106.5,9) {$-2$};
  \node[elem] (s3) at (113.5,9) {300};
  \node[addr, anchor=north west] at (96,6.3) {0x9c40};
  \node[lab, anchor=east] at (95.6,9) {\iflangja{ヒープ}{heap}};
  \draw[->] ([xshift=-4mm]sv.south) -- (s1.north);
  % ---------------- operations
  % Japanese: below the arrows; above them the wider "アンマーシャル"
  % runs into the buffer offset "16".
  \draw[->, thick, ln-accent] (31,20) -- (44.5,20);
  \node[op, anchor=\iflangja{north}{south}] at (37.8,\iflangja{19.2}{20.8}) {\iflangja{マーシャル}{marshal}};
  \draw[->, thick, ln-accent] (75.5,20) -- (87,20);
  \node[op, anchor=\iflangja{north}{south}] at (\iflangja{82.3}{81.3},\iflangja{19.2}{20.8}) {\iflangja{アンマーシャル}{unmarshal}};
\end{tikzpicture}

  \caption{Marshalling and unmarshalling the arguments of
    \texttt{sum(3, v)}.  The buffer contains the values of the arguments,
    including the array to which the pointer refers, but not the pointer
    itself; the numbers above it are byte offsets.}
  \label{fig:l13-marshal}
\end{figure}

The message that the runtime sends also identifies the procedure that is
being called; \cref{sec:l13-encoding} shows the header that Sun RPC places
in front of the encoded arguments.

Marshalling and unmarshalling code is not written by hand.  The compiler
of the IDL generates it from the interface, together with the stubs, as
\cref{sec:l13-rpcgen} shows for Sun RPC.

\section{Binding and registries}
\label{sec:l13-binding}

Before a client can call the procedures of a server,\source{P4L1, binding
and registry; Birrell and Nelson \cite{birrell1984}.} it must determine
which server to use and how to reach it.

\begin{definition}[Binding]
  In \term{binding}[バインディング] a client determines which server it will
  communicate with, identified by the name and the version of the service,
  and how it will connect to that server, given by the network address,
  the transport protocol and the port, and it establishes the
  communication.
\end{definition}

\begin{definition}[Registry]
  A \term{registry}[レジストリ] is a database of available services.
  Servers register their services in it, and clients search it by service
  name to find a server that offers the service and the details needed to
  contact that server.
\end{definition}

A registry can be organized in two ways.
\begin{description}
  \item[Distributed registry] A registry service, possibly replicated,
    with which any RPC service can register wherever it runs.  Clients find
    services without knowing their location.\margin{Birrell and Nelson used
    Grapevine, a distributed database, for binding \cite{birrell1984}.}
  \item[Machine-specific registry] Every machine runs a registry for the
    services on that machine.  A client must already know the address of
    the machine, and the registry provides the port number at which the
    service can be reached.
\end{description}
A registry also needs a naming protocol.  The simplest requires the service
name of the client to match the registered name exactly; a more elaborate
one could also consider synonyms, so that a client looking for a service
named \texttt{sum} would also find one registered as \texttt{add} or
\texttt{summation}.

\section{Pointers as arguments}
\label{sec:l13-pointers}

A local procedure such as \texttt{sum} of \cref{ex:l13-marshal} can
receive a pointer and dereference it, because caller and procedure share an
address space.\source{P4L1, pointers in RPCs; OSTEP ch.~48.}  In an RPC,
the pointer refers to the address space of the client.  On the server the
same address is meaningless: it may be unmapped or hold unrelated data.
An RPC system solves the problem in one of two ways.
\begin{description}
  \item[Disallow pointers] Remote interfaces may not contain pointer
    arguments.
  \item[Serialize the pointed-to data] The client stub copies the data
    structure to which the pointer refers into the send buffer, and the
    server stub recreates it in the address space of the server and passes
    a pointer to this copy, as in \cref{fig:l13-marshal}.  For linked data
    structures, the stubs follow the pointers recursively.
\end{description}
With serialized pointers the server operates on a copy: changes that it
makes to the pointed-to data are not visible to the client unless the
interface returns the data as part of the result.

\section{Partial failures}
\label{sec:l13-failures}

A local procedure call either returns or fails together with the process
that made it:\source{P4L1, handling partial failures, RPC design choice
summary; OSTEP ch.~48.} caller and procedure share the fate of a single
process on a single machine.  A distributed system can fail in parts.

\begin{definition}[Partial failure]
  A \term{partial failure}[部分障害] is a failure of some components of a
  distributed system, such as a machine, a process, a network link or an
  individual message, while the other components continue to operate.
\end{definition}

When an RPC does not return, the cause may be that the server machine is
down, that the server process has crashed or is overloaded, that the network
is down, or that the request or the reply has been lost.  The client
observes the same thing in all these cases---no reply---and cannot tell
which component has failed (\cref{fig:l13-failures}).

\begin{figure}[htbp]
  \centering
  % Partial failures seen by an RPC client: (a) request lost, (b) server
% crash, (c) reply lost followed by a retry that executes the call twice.
% Time runs downward.  Usage: % Partial failures seen by an RPC client: (a) request lost, (b) server
% crash, (c) reply lost followed by a retry that executes the call twice.
% Time runs downward.  Usage: % Partial failures seen by an RPC client: (a) request lost, (b) server
% crash, (c) reply lost followed by a retry that executes the call twice.
% Time runs downward.  Usage: \input{figures/l13-failures}  (width 120 mm)
\begin{tikzpicture}[x=1mm, y=0.85mm, font=\scriptsize\sffamily,
  >={Stealth[length=1.5mm]},
  life/.style={draw=ln-rule, thick},
  blocked/.style={draw=ln-accent, fill=ln-accent!25},
  exec/.style={draw, fill=ln-muted!45},
  msg/.style={->},
  hdr/.style={font=\scriptsize\sffamily\bfseries, text=ln-muted, anchor=base},
  lab/.style={font=\scriptsize\sffamily, text=ln-muted},
  pcap/.style={font=\footnotesize\sffamily, anchor=north}]
  \newcommand{\lnLost}[2]{%
    \draw[ln-caution, thick] (#1-1.1,#2-1.1) -- (#1+1.1,#2+1.1)
      (#1-1.1,#2+1.1) -- (#1+1.1,#2-1.1);}
  \newcommand{\lnTimeout}[1]{%
    \draw[ln-caution, thick] (1.5,#1) -- (6.5,#1);
    \node[lab, text=ln-caution, anchor=south west, inner sep=0.5pt]
      at (6.8,#1) {\iflangja{タイムアウト}{timeout}};}
  \foreach \xs/\pc in {0/{\iflangja{(a) 要求の消失}{(a) request lost}},
                        42/{\iflangja{(b) サーバの障害}{(b) server crash}},
                        84/{\iflangja{(c) 応答の消失}{(c) reply lost}}} {
    \begin{scope}[xshift=\xs mm]
      \node[hdr] at (4,45) {\iflangja{クライアント}{client}};
      \node[hdr] at (32,45) {\iflangja{サーバ}{server}};
      \node[pcap] at (18,-6) {\pc};
    \end{scope}
  }
  % ================= (a) request lost
  \draw[life] (4,42) -- (4,-4);
  \draw[life] (32,42) -- (32,-4);
  \draw[blocked] (3,36) rectangle (5,14);
  \draw[msg] (5,36) -- (17,32.2);
  \lnLost{18}{31.9}
  \lnTimeout{14}
  % ================= (b) server crash
  \begin{scope}[xshift=42mm]
  \draw[life] (4,42) -- (4,-4);
  \draw[life] (32,42) -- (32,27);
  \draw[life, densely dashed] (32,27) -- (32,-4);
  \draw[blocked] (3,36) rectangle (5,14);
  \draw[msg] (5,36) -- (31,30);
  \draw[exec] (31,30) rectangle (33,27.5);
  \lnLost{32}{26}
  \node[lab, anchor=east] at (30.5,24.5) {\iflangja{クラッシュ}{crash}};
  \lnTimeout{14}
  \end{scope}
  % ================= (c) reply lost, retry
  \begin{scope}[xshift=84mm]
  \draw[life] (4,42) -- (4,-4);
  \draw[life] (32,42) -- (32,-4);
  \draw[blocked] (3,36) rectangle (5,-1);
  \draw[msg] (5,36) -- (31,30);
  \draw[exec] (31,30) rectangle (33,26);
  \draw[msg] (31,26) -- (19,22.3);
  \lnLost{18}{22}
  \lnTimeout{14}
  \draw[msg] (5,12) -- (31,6);
  \node[lab, anchor=south, rotate=-11] at (18,9.6) {\iflangja{再送}{retry}};
  \draw[exec] (31,6) rectangle (33,2);
  \draw[msg] (31,2) -- (5,-1);
  \node[lab, anchor=east] at (30.5,4) {\iflangja{再実行}{runs again}};
  \end{scope}
\end{tikzpicture}
  (width 120 mm)
\begin{tikzpicture}[x=1mm, y=0.85mm, font=\scriptsize\sffamily,
  >={Stealth[length=1.5mm]},
  life/.style={draw=ln-rule, thick},
  blocked/.style={draw=ln-accent, fill=ln-accent!25},
  exec/.style={draw, fill=ln-muted!45},
  msg/.style={->},
  hdr/.style={font=\scriptsize\sffamily\bfseries, text=ln-muted, anchor=base},
  lab/.style={font=\scriptsize\sffamily, text=ln-muted},
  pcap/.style={font=\footnotesize\sffamily, anchor=north}]
  \newcommand{\lnLost}[2]{%
    \draw[ln-caution, thick] (#1-1.1,#2-1.1) -- (#1+1.1,#2+1.1)
      (#1-1.1,#2+1.1) -- (#1+1.1,#2-1.1);}
  \newcommand{\lnTimeout}[1]{%
    \draw[ln-caution, thick] (1.5,#1) -- (6.5,#1);
    \node[lab, text=ln-caution, anchor=south west, inner sep=0.5pt]
      at (6.8,#1) {\iflangja{タイムアウト}{timeout}};}
  \foreach \xs/\pc in {0/{\iflangja{(a) 要求の消失}{(a) request lost}},
                        42/{\iflangja{(b) サーバの障害}{(b) server crash}},
                        84/{\iflangja{(c) 応答の消失}{(c) reply lost}}} {
    \begin{scope}[xshift=\xs mm]
      \node[hdr] at (4,45) {\iflangja{クライアント}{client}};
      \node[hdr] at (32,45) {\iflangja{サーバ}{server}};
      \node[pcap] at (18,-6) {\pc};
    \end{scope}
  }
  % ================= (a) request lost
  \draw[life] (4,42) -- (4,-4);
  \draw[life] (32,42) -- (32,-4);
  \draw[blocked] (3,36) rectangle (5,14);
  \draw[msg] (5,36) -- (17,32.2);
  \lnLost{18}{31.9}
  \lnTimeout{14}
  % ================= (b) server crash
  \begin{scope}[xshift=42mm]
  \draw[life] (4,42) -- (4,-4);
  \draw[life] (32,42) -- (32,27);
  \draw[life, densely dashed] (32,27) -- (32,-4);
  \draw[blocked] (3,36) rectangle (5,14);
  \draw[msg] (5,36) -- (31,30);
  \draw[exec] (31,30) rectangle (33,27.5);
  \lnLost{32}{26}
  \node[lab, anchor=east] at (30.5,24.5) {\iflangja{クラッシュ}{crash}};
  \lnTimeout{14}
  \end{scope}
  % ================= (c) reply lost, retry
  \begin{scope}[xshift=84mm]
  \draw[life] (4,42) -- (4,-4);
  \draw[life] (32,42) -- (32,-4);
  \draw[blocked] (3,36) rectangle (5,-1);
  \draw[msg] (5,36) -- (31,30);
  \draw[exec] (31,30) rectangle (33,26);
  \draw[msg] (31,26) -- (19,22.3);
  \lnLost{18}{22}
  \lnTimeout{14}
  \draw[msg] (5,12) -- (31,6);
  \node[lab, anchor=south, rotate=-11] at (18,9.6) {\iflangja{再送}{retry}};
  \draw[exec] (31,6) rectangle (33,2);
  \draw[msg] (31,2) -- (5,-1);
  \node[lab, anchor=east] at (30.5,4) {\iflangja{再実行}{runs again}};
  \end{scope}
\end{tikzpicture}
  (width 120 mm)
\begin{tikzpicture}[x=1mm, y=0.85mm, font=\scriptsize\sffamily,
  >={Stealth[length=1.5mm]},
  life/.style={draw=ln-rule, thick},
  blocked/.style={draw=ln-accent, fill=ln-accent!25},
  exec/.style={draw, fill=ln-muted!45},
  msg/.style={->},
  hdr/.style={font=\scriptsize\sffamily\bfseries, text=ln-muted, anchor=base},
  lab/.style={font=\scriptsize\sffamily, text=ln-muted},
  pcap/.style={font=\footnotesize\sffamily, anchor=north}]
  \newcommand{\lnLost}[2]{%
    \draw[ln-caution, thick] (#1-1.1,#2-1.1) -- (#1+1.1,#2+1.1)
      (#1-1.1,#2+1.1) -- (#1+1.1,#2-1.1);}
  \newcommand{\lnTimeout}[1]{%
    \draw[ln-caution, thick] (1.5,#1) -- (6.5,#1);
    \node[lab, text=ln-caution, anchor=south west, inner sep=0.5pt]
      at (6.8,#1) {\iflangja{タイムアウト}{timeout}};}
  \foreach \xs/\pc in {0/{\iflangja{(a) 要求の消失}{(a) request lost}},
                        42/{\iflangja{(b) サーバの障害}{(b) server crash}},
                        84/{\iflangja{(c) 応答の消失}{(c) reply lost}}} {
    \begin{scope}[xshift=\xs mm]
      \node[hdr] at (4,45) {\iflangja{クライアント}{client}};
      \node[hdr] at (32,45) {\iflangja{サーバ}{server}};
      \node[pcap] at (18,-6) {\pc};
    \end{scope}
  }
  % ================= (a) request lost
  \draw[life] (4,42) -- (4,-4);
  \draw[life] (32,42) -- (32,-4);
  \draw[blocked] (3,36) rectangle (5,14);
  \draw[msg] (5,36) -- (17,32.2);
  \lnLost{18}{31.9}
  \lnTimeout{14}
  % ================= (b) server crash
  \begin{scope}[xshift=42mm]
  \draw[life] (4,42) -- (4,-4);
  \draw[life] (32,42) -- (32,27);
  \draw[life, densely dashed] (32,27) -- (32,-4);
  \draw[blocked] (3,36) rectangle (5,14);
  \draw[msg] (5,36) -- (31,30);
  \draw[exec] (31,30) rectangle (33,27.5);
  \lnLost{32}{26}
  \node[lab, anchor=east] at (30.5,24.5) {\iflangja{クラッシュ}{crash}};
  \lnTimeout{14}
  \end{scope}
  % ================= (c) reply lost, retry
  \begin{scope}[xshift=84mm]
  \draw[life] (4,42) -- (4,-4);
  \draw[life] (32,42) -- (32,-4);
  \draw[blocked] (3,36) rectangle (5,-1);
  \draw[msg] (5,36) -- (31,30);
  \draw[exec] (31,30) rectangle (33,26);
  \draw[msg] (31,26) -- (19,22.3);
  \lnLost{18}{22}
  \lnTimeout{14}
  \draw[msg] (5,12) -- (31,6);
  \node[lab, anchor=south, rotate=-11] at (18,9.6) {\iflangja{再送}{retry}};
  \draw[exec] (31,6) rectangle (33,2);
  \draw[msg] (31,2) -- (5,-1);
  \node[lab, anchor=east] at (30.5,4) {\iflangja{再実行}{runs again}};
  \end{scope}
\end{tikzpicture}

  \caption{Three partial failures that the client cannot distinguish.  The
    shaded bar marks the time during which the client is blocked, grey
    boxes the execution of the procedure, and a cross a lost message.
    (a)~The request is lost; (b)~the
    server crashes while executing the call; (c)~the reply is lost, and the
    retry executes the procedure a second time.}
  \label{fig:l13-failures}
\end{figure}

An RPC system can bound the waiting with a \emph{timeout} and retransmit
the request when it expires, but timeouts and retries give no guarantees.
Retries cannot succeed while the server is down, and if the reply was lost,
a retry executes the procedure again, as in \cref{fig:l13-failures}c.
Repeated execution is harmless only for \emph{idempotent} procedures,
whose repeated execution has the same effect as a single one.

\begin{example}
  A procedure that appends a record to a file is not idempotent: if the
  reply to a call that appends a \SI{200}{\byte} record is lost, the retry
  appends the record a second time and the file grows by \SI{400}{\byte}.
  A procedure that writes \SI{200}{\byte} at a given offset, such as offset
  \num{4096}, is idempotent: executing it twice leaves the same contents.
\end{example}

RPC systems therefore provide a special error notification---an error
code, a signal or an exception---that serves as a catch-all for the ways in
which a call can fail partially.  The caller must be prepared to receive it
on every call, although it does not say which component has failed or
whether the procedure has executed.\margin{Birrell and Nelson: exactly-once
execution if the call returns, otherwise an exception \cite{birrell1984}.}

\begin{keypoint}[RPC design choices]
  Every RPC system must decide
  \begin{itemize}
    \item how to find the server: binding, through a registry;
    \item how to talk to the server and package the data: an IDL and
      marshalling;
    \item how to treat pointer arguments: disallow them or serialize the
      pointed-to data;
    \item how to handle partial failures: with special error notifications.
  \end{itemize}
\end{keypoint}

\section{Sun RPC}
\label{sec:l13-sunrpc}

Sun Microsystems developed \term{Sun RPC}[Sun RPC] in the 1980s for Unix
systems; its best-known application is the Network File System
(Lesson~14).\source{P4L1, what is SunRPC, SunRPC overview, SunRPC XDR
example; RFC 5531 \cite{rfc5531}; RFC 4506 \cite{rfc4506}.}  It is now
available on most platforms and has been standardized as ONC~RPC
(Open Network Computing RPC) \cite{rfc5531}.\index{ONC RPC|see{Sun RPC}}
Sun RPC makes the four design choices summarized at the end of
\cref{sec:l13-failures} as follows.
\begin{description}
  \item[Binding] A registry daemon runs on every machine that offers RPC
    services (\cref{sec:l13-sun-binding}).
  \item[IDL] XDR serves both for specifying interfaces and for encoding data
    on the wire.
  \item[Pointers] Pointers are allowed, and the pointed-to data are
    serialized.
  \item[Failures] The runtime retries calls, using timeouts, and returns as
    much information as possible about the cause of a failure.
\end{description}

\begin{definition}[External data representation]
  The language-agnostic IDL of Sun RPC, together with the standard binary
  encoding of the data that it describes, is called
  \term{external data representation}[外部データ表現] (XDR)
  \cite{rfc4506}.
\end{definition}
\index{XDR|see{external data representation}}

Client and server interact through procedure calls.  The interface of a
service is specified in XDR in a \texttt{.x} file, from which
\texttt{rpcgen} generates stubs for a specific language, C in the following
(\cref{sec:l13-rpcgen}).  A server registers with the registry daemon of its
machine, which records the number of the service, its version, the
transport protocols and the port.  Binding creates a \emph{client handle},
which the client passes to every call and in which the runtime keeps the
state of the client's RPCs.  Client and server may run on the same machine
or on different machines.  Current implementations descend from TI-RPC
(transport-independent RPC), whose library does not depend on a particular
transport protocol.

\subsection{Interface specification}

The following XDR specification, \texttt{sum.x}, describes a service that
returns the sum of up to 64 integers.

\begin{lstlisting}[language=C, numbers=none]
const MAX_VALUES = 64;

struct sum_in {
  int values<MAX_VALUES>;    /* variable-length array */
};
struct sum_out {
  int total;
};

program SUM_PROG {                     /* service */
  version SUM_VERS {
    sum_out SUM_PROC(sum_in) = 1;      /* procedure 1 */
  } = 1;                               /* version 1 */
} = 0x20000123;                        /* program number */
\end{lstlisting}

The data types \texttt{sum\_in} and \texttt{sum\_out} describe the argument
and the result.  A procedure takes a single argument and returns a single
result, so several values are grouped in a structure.  The service itself
is a \emph{program} that contains one or more versions, each of which
contains procedures; this part of the \texttt{.x} language is the RPC
language, an extension of XDR defined together with the protocol
\cite{rfc5531}.  Program, version and procedure are identified by numbers,
and a request names all three.  The 32-bit program number identifies the
service uniquely.  A client that sets
\texttt{values} to the three integers 7, $-2$ and 300 and calls
\texttt{SUM\_PROC} receives a \texttt{sum\_out} with \texttt{total} $=305$.

\section{Compiling an interface with rpcgen}
\label{sec:l13-rpcgen}

The stubs of Sun RPC are generated by \term{rpcgen}[rpcgen], the compiler
that translates an XDR specification into C code.\source{P4L1, compiling
XDR, summarizing XDR compilation; OSTEP ch.~48.}  The command
\texttt{rpcgen -C sum.x}, where \texttt{-C} selects ANSI~C, produces four
files (\cref{fig:l13-rpcgen}).
\begin{description}
  \item[\texttt{sum.h}] The C declarations of the data types and of the
    stub procedures, and constants such as \texttt{SUM\_PROG},
    \texttt{SUM\_VERS} and \texttt{SUM\_PROC}.  All other files include it.
  \item[\texttt{sum\_clnt.c}] The client stub \texttt{sum\_proc\_1}, named
    after the procedure and the version.  It marshals the argument, sends
    the request through the runtime, waits for the reply and unmarshals the
    result.
  \item[\texttt{sum\_svc.c}] The server stub and a skeleton of the server
    program.  Its \texttt{main} function creates the transports, registers
    the program and version with the registry daemon, and enters a loop
    that waits for requests.  For each request, the dispatch procedure
    \texttt{sum\_prog\_1} determines the procedure number, unmarshals the
    argument, calls \texttt{sum\_proc\_1\_svc}, and marshals and sends the
    result.
  \item[\texttt{sum\_xdr.c}] The XDR routines that marshal and unmarshal
    \texttt{sum\_in} and \texttt{sum\_out}, used by both sides
    (\cref{sec:l13-xdr-types}).
\end{description}

\begin{figure}[htbp]
  \centering
  % Building a Sun RPC client and server with rpcgen from sum.x.
% White: written by the developer; shaded: generated by rpcgen.
% Usage: % Building a Sun RPC client and server with rpcgen from sum.x.
% White: written by the developer; shaded: generated by rpcgen.
% Usage: % Building a Sun RPC client and server with rpcgen from sum.x.
% White: written by the developer; shaded: generated by rpcgen.
% Usage: \input{figures/l13-rpcgen}   (width ~120 mm)
\begin{tikzpicture}[x=1mm, y=1mm, font=\footnotesize\sffamily,
  >={Stealth[length=1.5mm]},
  file/.style={draw, fill=white, align=center, minimum width=18.5mm,
               minimum height=9mm, inner sep=0.5pt},
  gen/.style={file, fill=ln-accent!15},
  tool/.style={draw, rounded corners=3pt, fill=ln-accent!40, align=center,
               minimum width=22mm, minimum height=6.5mm},
  exe/.style={draw, thick, rounded corners=2pt, fill=ln-box, align=center,
              minimum width=30mm, minimum height=7mm},
  fn/.style={font=\scriptsize\ttfamily},
  what/.style={font=\tiny\sffamily, text=ln-muted},
  lab/.style={font=\scriptsize\sffamily, text=ln-muted}]
  % English descriptions at \scriptsize: \tiny is unreadable on screen, and
  % the widest ("types, constants", 16.4 mm) still fits the 18.5 mm box.
  \newcommand{\lnFile}[2]{{\scriptsize\ttfamily #1}\\[-0.4ex]{\iflangja{\tiny}{\scriptsize}\sffamily\color{ln-muted}#2}}
  % interface and compiler
  \node[file] (x) at (60,52) {\lnFile{sum.x}{\iflangja{XDR インタフェース}{XDR interface}}};
  \node[tool] (rpcgen) at (60,40) {\texttt{rpcgen}};
  \draw[->] (x) -- (rpcgen);
  % source files
  \node[file] (cli) at (10,24) {\lnFile{client.c}{\iflangja{呼び出し元}{caller}}};
  \node[gen]  (cln) at (30,24) {\lnFile{sum\_clnt.c}{\iflangja{クライアントスタブ}{client stub}}};
  \node[gen]  (h)   at (50,24) {\lnFile{sum.h}{\iflangja{型と定数}{types, constants}}};
  \node[gen]  (xdr) at (70,24) {\lnFile{sum\_xdr.c}{\iflangja{XDR ルーチン}{XDR routines}}};
  \node[gen]  (svc) at (90,24) {\lnFile{sum\_svc.c}{\iflangja{サーバスタブ}{server stub}}};
  \node[file] (srv) at (110,24) {\lnFile{server.c}{\iflangja{手続きの実装}{procedure}}};
  \foreach \f in {cln, h, xdr, svc} \draw[->] (rpcgen.south) -- (\f.north);
  % legend
  \node[file, minimum width=5mm, minimum height=3.5mm] (l1) at (3,53) {};
  \node[lab, anchor=west] at (6,53) {\iflangja{開発者が記述}{written by the developer}};
  \node[gen, minimum width=5mm, minimum height=3.5mm] (l2) at (3,48) {};
  \node[lab, anchor=west] at (6,48) {\iflangja{\texttt{rpcgen} が生成}{generated by \texttt{rpcgen}}};
  % executables
  \node[exe] (ce) at (30,5) {\iflangja{クライアントプログラム}{client program}};
  \node[exe] (se) at (90,5) {\iflangja{サーバプログラム}{server program}};
  \foreach \f in {cli, cln, xdr} \draw[->] (\f.south) -- (ce.north);
  \foreach \f in {xdr, svc, srv} \draw[->] (\f.south) -- (se.north);
  \node[lab] at (60,12.5) {\iflangja{コンパイル・リンク}{compile and link}};
\end{tikzpicture}
   (width ~120 mm)
\begin{tikzpicture}[x=1mm, y=1mm, font=\footnotesize\sffamily,
  >={Stealth[length=1.5mm]},
  file/.style={draw, fill=white, align=center, minimum width=18.5mm,
               minimum height=9mm, inner sep=0.5pt},
  gen/.style={file, fill=ln-accent!15},
  tool/.style={draw, rounded corners=3pt, fill=ln-accent!40, align=center,
               minimum width=22mm, minimum height=6.5mm},
  exe/.style={draw, thick, rounded corners=2pt, fill=ln-box, align=center,
              minimum width=30mm, minimum height=7mm},
  fn/.style={font=\scriptsize\ttfamily},
  what/.style={font=\tiny\sffamily, text=ln-muted},
  lab/.style={font=\scriptsize\sffamily, text=ln-muted}]
  % English descriptions at \scriptsize: \tiny is unreadable on screen, and
  % the widest ("types, constants", 16.4 mm) still fits the 18.5 mm box.
  \newcommand{\lnFile}[2]{{\scriptsize\ttfamily #1}\\[-0.4ex]{\iflangja{\tiny}{\scriptsize}\sffamily\color{ln-muted}#2}}
  % interface and compiler
  \node[file] (x) at (60,52) {\lnFile{sum.x}{\iflangja{XDR インタフェース}{XDR interface}}};
  \node[tool] (rpcgen) at (60,40) {\texttt{rpcgen}};
  \draw[->] (x) -- (rpcgen);
  % source files
  \node[file] (cli) at (10,24) {\lnFile{client.c}{\iflangja{呼び出し元}{caller}}};
  \node[gen]  (cln) at (30,24) {\lnFile{sum\_clnt.c}{\iflangja{クライアントスタブ}{client stub}}};
  \node[gen]  (h)   at (50,24) {\lnFile{sum.h}{\iflangja{型と定数}{types, constants}}};
  \node[gen]  (xdr) at (70,24) {\lnFile{sum\_xdr.c}{\iflangja{XDR ルーチン}{XDR routines}}};
  \node[gen]  (svc) at (90,24) {\lnFile{sum\_svc.c}{\iflangja{サーバスタブ}{server stub}}};
  \node[file] (srv) at (110,24) {\lnFile{server.c}{\iflangja{手続きの実装}{procedure}}};
  \foreach \f in {cln, h, xdr, svc} \draw[->] (rpcgen.south) -- (\f.north);
  % legend
  \node[file, minimum width=5mm, minimum height=3.5mm] (l1) at (3,53) {};
  \node[lab, anchor=west] at (6,53) {\iflangja{開発者が記述}{written by the developer}};
  \node[gen, minimum width=5mm, minimum height=3.5mm] (l2) at (3,48) {};
  \node[lab, anchor=west] at (6,48) {\iflangja{\texttt{rpcgen} が生成}{generated by \texttt{rpcgen}}};
  % executables
  \node[exe] (ce) at (30,5) {\iflangja{クライアントプログラム}{client program}};
  \node[exe] (se) at (90,5) {\iflangja{サーバプログラム}{server program}};
  \foreach \f in {cli, cln, xdr} \draw[->] (\f.south) -- (ce.north);
  \foreach \f in {xdr, svc, srv} \draw[->] (\f.south) -- (se.north);
  \node[lab] at (60,12.5) {\iflangja{コンパイル・リンク}{compile and link}};
\end{tikzpicture}
   (width ~120 mm)
\begin{tikzpicture}[x=1mm, y=1mm, font=\footnotesize\sffamily,
  >={Stealth[length=1.5mm]},
  file/.style={draw, fill=white, align=center, minimum width=18.5mm,
               minimum height=9mm, inner sep=0.5pt},
  gen/.style={file, fill=ln-accent!15},
  tool/.style={draw, rounded corners=3pt, fill=ln-accent!40, align=center,
               minimum width=22mm, minimum height=6.5mm},
  exe/.style={draw, thick, rounded corners=2pt, fill=ln-box, align=center,
              minimum width=30mm, minimum height=7mm},
  fn/.style={font=\scriptsize\ttfamily},
  what/.style={font=\tiny\sffamily, text=ln-muted},
  lab/.style={font=\scriptsize\sffamily, text=ln-muted}]
  % English descriptions at \scriptsize: \tiny is unreadable on screen, and
  % the widest ("types, constants", 16.4 mm) still fits the 18.5 mm box.
  \newcommand{\lnFile}[2]{{\scriptsize\ttfamily #1}\\[-0.4ex]{\iflangja{\tiny}{\scriptsize}\sffamily\color{ln-muted}#2}}
  % interface and compiler
  \node[file] (x) at (60,52) {\lnFile{sum.x}{\iflangja{XDR インタフェース}{XDR interface}}};
  \node[tool] (rpcgen) at (60,40) {\texttt{rpcgen}};
  \draw[->] (x) -- (rpcgen);
  % source files
  \node[file] (cli) at (10,24) {\lnFile{client.c}{\iflangja{呼び出し元}{caller}}};
  \node[gen]  (cln) at (30,24) {\lnFile{sum\_clnt.c}{\iflangja{クライアントスタブ}{client stub}}};
  \node[gen]  (h)   at (50,24) {\lnFile{sum.h}{\iflangja{型と定数}{types, constants}}};
  \node[gen]  (xdr) at (70,24) {\lnFile{sum\_xdr.c}{\iflangja{XDR ルーチン}{XDR routines}}};
  \node[gen]  (svc) at (90,24) {\lnFile{sum\_svc.c}{\iflangja{サーバスタブ}{server stub}}};
  \node[file] (srv) at (110,24) {\lnFile{server.c}{\iflangja{手続きの実装}{procedure}}};
  \foreach \f in {cln, h, xdr, svc} \draw[->] (rpcgen.south) -- (\f.north);
  % legend
  \node[file, minimum width=5mm, minimum height=3.5mm] (l1) at (3,53) {};
  \node[lab, anchor=west] at (6,53) {\iflangja{開発者が記述}{written by the developer}};
  \node[gen, minimum width=5mm, minimum height=3.5mm] (l2) at (3,48) {};
  \node[lab, anchor=west] at (6,48) {\iflangja{\texttt{rpcgen} が生成}{generated by \texttt{rpcgen}}};
  % executables
  \node[exe] (ce) at (30,5) {\iflangja{クライアントプログラム}{client program}};
  \node[exe] (se) at (90,5) {\iflangja{サーバプログラム}{server program}};
  \foreach \f in {cli, cln, xdr} \draw[->] (\f.south) -- (ce.north);
  \foreach \f in {xdr, svc, srv} \draw[->] (\f.south) -- (se.north);
  \node[lab] at (60,12.5) {\iflangja{コンパイル・リンク}{compile and link}};
\end{tikzpicture}

  \caption{Building a Sun RPC client and server.  The developer writes the
    interface, the client program and the implementation of the procedure;
    \texttt{rpcgen} generates the rest.  \texttt{sum.h} is included by all
    C files.}
  \label{fig:l13-rpcgen}
\end{figure}

The developer writes two pieces of code.  On the server side it is the
implementation of the procedure, \texttt{sum\_proc\_1\_svc}, which follows;
on the client side it is the program that includes \texttt{sum.h} and calls
\texttt{sum\_proc\_1} (\cref{sec:l13-sun-binding}).  Each is compiled and
linked with the generated files.  The RPC runtime library does the rest: the interaction
with the operating system, such as the use of sockets, and the management
of the communication.

\begin{lstlisting}[language=C, numbers=none]
#include "sum.h"

sum_out *
sum_proc_1_svc(sum_in *in, struct svc_req *req)
{
  static sum_out out;       /* must outlive the call */
  u_int i;

  out.total = 0;
  for (i = 0; i < in->values.values_len; i++)
    out.total += in->values.values_val[i];
  return &out;              /* marshalled by the stub */
}
\end{lstlisting}

The result is returned by pointer, and it must remain valid after the
procedure has returned, until the server stub has marshalled it, so it is
kept in a static variable.  The generated client stub likewise returns a
pointer to a static variable, which the next call overwrites.  This makes
the code generated by \texttt{rpcgen -C} unsafe for multithreaded clients
and servers.  With the option \texttt{-M}, \texttt{rpcgen} generates
thread-safe code, in which the caller of a stub supplies the memory for the
result.  The option does not make the server multithreaded: the generated
\texttt{main} still handles one request at a time.  On Solaris, the
additional option \texttt{-A} lets the runtime of the server create threads
for requests automatically.

\section{Registry and binding in Sun RPC}
\label{sec:l13-sun-binding}

The registry of Sun RPC is machine-specific:\source{P4L1, SunRPC registry
and binding; RFC 1833 \cite{rfc1833}.} each machine that offers RPC
services runs a registry daemon of its own, which clients contact at a
well-known port.

\begin{definition}[rpcbind]
  The registry daemon of Sun RPC is \term{rpcbind}[rpcbind].  It listens
  on port 111 for both TCP and UDP and maps a program number, version
  number and transport protocol to the address at which the corresponding
  server listens.
\end{definition}
\index{portmapper|see{rpcbind}}

The daemon \texttt{rpcbind} succeeded the older \emph{portmapper} daemon,
whose name still appears in the output below.  Because port 111 is a
privileged port, the daemon must be started with superuser privileges.  When a server starts, the generated
\texttt{main} registers each version of the program for each transport
(\cref{fig:l13-rpcbind}, step~1).  The registrations of a machine can be
listed with \texttt{rpcinfo -p}:

\begin{lstlisting}[numbers=none]
$ rpcinfo -p
   program vers proto   port  service
    100000    4   tcp    111  portmapper
    100000    2   udp    111  portmapper
    100003    3   tcp   2049  nfs
 536871203    1   tcp  40123
\end{lstlisting}

Each line gives the program number, the version, the transport protocol
and the port.  The daemon itself appears as program \texttt{100000} on
port 111, and the NFS server as program \texttt{100003}.  The last line is
the \texttt{sum} server, whose program number \texttt{0x20000123} is
printed in decimal.

\begin{figure}[htbp]
  \centering
  % Registration and binding in Sun RPC through the rpcbind daemon.
% 1 register, 2 query, 3 port number, 4 calls sent to the server's port.
% Usage: % Registration and binding in Sun RPC through the rpcbind daemon.
% 1 register, 2 query, 3 port number, 4 calls sent to the server's port.
% Usage: % Registration and binding in Sun RPC through the rpcbind daemon.
% 1 register, 2 query, 3 port number, 4 calls sent to the server's port.
% Usage: \input{figures/l13-rpcbind}   (width ~120 mm)
\begin{tikzpicture}[x=1mm, y=1mm, font=\footnotesize\sffamily,
  >={Stealth[length=1.6mm]},
  machine/.style={draw=ln-rule, rounded corners=2pt, fill=ln-box},
  proc/.style={draw, fill=white, align=center, inner sep=1pt},
  daemon/.style={proc, fill=ln-accent!20},
  stepn/.style={circle, draw, fill=white, inner sep=0pt,
                minimum size=4.2mm, font=\scriptsize\sffamily\bfseries},
  hdr/.style={font=\scriptsize\sffamily\bfseries, text=ln-muted},
  msg/.style={font=\scriptsize\sffamily, inner sep=1pt}]
  \draw[machine] (0,0) rectangle (40,34);
  \draw[machine] (64,0) rectangle (120,34);
  \node[hdr] at (20,31) {\iflangja{クライアントマシン}{client machine}};
  \node[hdr] at (92,31) {\iflangja{サーバマシン}{server machine}};
  \node[proc, minimum width=32mm, minimum height=22mm] (cli) at (20,15)
    {\iflangja{クライアント}{client}\\[0.5ex]\texttt{clnt\_create()}\\[0.5ex]
     \texttt{sum\_proc\_1()}};
  \node[daemon, minimum width=40mm, minimum height=9mm] (rb) at (92,22)
    {\texttt{rpcbind}\\[-0.3ex]\iflangja{ポート}{port} 111};
  \node[proc, minimum width=40mm, minimum height=9mm] (srv) at (92,7)
    {\iflangja{\texttt{sum} サーバ}{\texttt{sum} server}\\[-0.3ex]\iflangja{ポート}{port} 40123};
  % 1 register
  \draw[->] ([xshift=12mm]srv.north) -- ([xshift=12mm]rb.south);
  \node[stepn] at (108.5,14.5) {1};
  % 2 query, 3 reply
  \draw[->] (36,24) -- (72,24) node[pos=0.6, above, msg]
    {\texttt{(SUM\_PROG, 1, tcp)}};
  \draw[<-] (36,20) -- (72,20) node[pos=0.6, below, msg]
    {\iflangja{ポート}{port} 40123};
  \node[stepn] at (40,27) {2};
  \node[stepn] at (40,17) {3};
  % 4 calls
  \draw[->, thick, ln-accent] (36,7) -- (72,7) node[pos=0.6, above, msg, text=black]
    {\iflangja{RPC 呼び出し}{RPC calls}};
  \node[stepn] at (40,3.5) {4};
\end{tikzpicture}
   (width ~120 mm)
\begin{tikzpicture}[x=1mm, y=1mm, font=\footnotesize\sffamily,
  >={Stealth[length=1.6mm]},
  machine/.style={draw=ln-rule, rounded corners=2pt, fill=ln-box},
  proc/.style={draw, fill=white, align=center, inner sep=1pt},
  daemon/.style={proc, fill=ln-accent!20},
  stepn/.style={circle, draw, fill=white, inner sep=0pt,
                minimum size=4.2mm, font=\scriptsize\sffamily\bfseries},
  hdr/.style={font=\scriptsize\sffamily\bfseries, text=ln-muted},
  msg/.style={font=\scriptsize\sffamily, inner sep=1pt}]
  \draw[machine] (0,0) rectangle (40,34);
  \draw[machine] (64,0) rectangle (120,34);
  \node[hdr] at (20,31) {\iflangja{クライアントマシン}{client machine}};
  \node[hdr] at (92,31) {\iflangja{サーバマシン}{server machine}};
  \node[proc, minimum width=32mm, minimum height=22mm] (cli) at (20,15)
    {\iflangja{クライアント}{client}\\[0.5ex]\texttt{clnt\_create()}\\[0.5ex]
     \texttt{sum\_proc\_1()}};
  \node[daemon, minimum width=40mm, minimum height=9mm] (rb) at (92,22)
    {\texttt{rpcbind}\\[-0.3ex]\iflangja{ポート}{port} 111};
  \node[proc, minimum width=40mm, minimum height=9mm] (srv) at (92,7)
    {\iflangja{\texttt{sum} サーバ}{\texttt{sum} server}\\[-0.3ex]\iflangja{ポート}{port} 40123};
  % 1 register
  \draw[->] ([xshift=12mm]srv.north) -- ([xshift=12mm]rb.south);
  \node[stepn] at (108.5,14.5) {1};
  % 2 query, 3 reply
  \draw[->] (36,24) -- (72,24) node[pos=0.6, above, msg]
    {\texttt{(SUM\_PROG, 1, tcp)}};
  \draw[<-] (36,20) -- (72,20) node[pos=0.6, below, msg]
    {\iflangja{ポート}{port} 40123};
  \node[stepn] at (40,27) {2};
  \node[stepn] at (40,17) {3};
  % 4 calls
  \draw[->, thick, ln-accent] (36,7) -- (72,7) node[pos=0.6, above, msg, text=black]
    {\iflangja{RPC 呼び出し}{RPC calls}};
  \node[stepn] at (40,3.5) {4};
\end{tikzpicture}
   (width ~120 mm)
\begin{tikzpicture}[x=1mm, y=1mm, font=\footnotesize\sffamily,
  >={Stealth[length=1.6mm]},
  machine/.style={draw=ln-rule, rounded corners=2pt, fill=ln-box},
  proc/.style={draw, fill=white, align=center, inner sep=1pt},
  daemon/.style={proc, fill=ln-accent!20},
  stepn/.style={circle, draw, fill=white, inner sep=0pt,
                minimum size=4.2mm, font=\scriptsize\sffamily\bfseries},
  hdr/.style={font=\scriptsize\sffamily\bfseries, text=ln-muted},
  msg/.style={font=\scriptsize\sffamily, inner sep=1pt}]
  \draw[machine] (0,0) rectangle (40,34);
  \draw[machine] (64,0) rectangle (120,34);
  \node[hdr] at (20,31) {\iflangja{クライアントマシン}{client machine}};
  \node[hdr] at (92,31) {\iflangja{サーバマシン}{server machine}};
  \node[proc, minimum width=32mm, minimum height=22mm] (cli) at (20,15)
    {\iflangja{クライアント}{client}\\[0.5ex]\texttt{clnt\_create()}\\[0.5ex]
     \texttt{sum\_proc\_1()}};
  \node[daemon, minimum width=40mm, minimum height=9mm] (rb) at (92,22)
    {\texttt{rpcbind}\\[-0.3ex]\iflangja{ポート}{port} 111};
  \node[proc, minimum width=40mm, minimum height=9mm] (srv) at (92,7)
    {\iflangja{\texttt{sum} サーバ}{\texttt{sum} server}\\[-0.3ex]\iflangja{ポート}{port} 40123};
  % 1 register
  \draw[->] ([xshift=12mm]srv.north) -- ([xshift=12mm]rb.south);
  \node[stepn] at (108.5,14.5) {1};
  % 2 query, 3 reply
  \draw[->] (36,24) -- (72,24) node[pos=0.6, above, msg]
    {\texttt{(SUM\_PROG, 1, tcp)}};
  \draw[<-] (36,20) -- (72,20) node[pos=0.6, below, msg]
    {\iflangja{ポート}{port} 40123};
  \node[stepn] at (40,27) {2};
  \node[stepn] at (40,17) {3};
  % 4 calls
  \draw[->, thick, ln-accent] (36,7) -- (72,7) node[pos=0.6, above, msg, text=black]
    {\iflangja{RPC 呼び出し}{RPC calls}};
  \node[stepn] at (40,3.5) {4};
\end{tikzpicture}

  \caption{Registration and binding in Sun RPC.  (1)~The server registers
    its program, version, transport and port with \texttt{rpcbind};
    (2,\,3)~\texttt{clnt\_create} asks \texttt{rpcbind} on the server
    machine for the port; (4)~the calls go directly to the server.}
  \label{fig:l13-rpcbind}
\end{figure}

A client binds with \texttt{clnt\_create}, whose arguments are the name of
the server machine, the program number, the version number and the
transport protocol.  \texttt{clnt\_create} asks \texttt{rpcbind} on that
machine for the port of the server (steps 2 and~3) and returns a pointer to
a client handle of type \texttt{CLIENT}, which the client passes to every
stub.  Besides the address of the server, the handle holds the
authentication information and the status of the most recent call, from
which error messages are derived.  The complete client follows.

\begin{lstlisting}[language=C, numbers=none]
#include <stdio.h>
#include <stdlib.h>
#include "sum.h"

int main(int argc, char *argv[])
{
  int v[3] = { 7, -2, 300 };
  sum_in in = { { 3, v } };   /* values_len, values_val */
  sum_out *out;
  CLIENT *clnt;

  if (argc != 2) {
    fprintf(stderr, "usage: %s host\n", argv[0]);
    exit(1);
  }
  /* bind: ask rpcbind on host argv[1] for the port */
  clnt = clnt_create(argv[1], SUM_PROG, SUM_VERS, "tcp");
  if (clnt == NULL) {
    clnt_pcreateerror(argv[1]);
    exit(1);
  }
  out = sum_proc_1(&in, clnt);         /* the RPC */
  if (out == NULL) {                   /* RPC failed */
    clnt_perror(clnt, argv[1]);
    exit(1);
  }
  printf("total = %d\n", out->total);
  clnt_destroy(clnt);
  return 0;
}
\end{lstlisting}

A stub that returns \texttt{NULL} reports a failed call.  The runtime has
recorded the reason in the handle---for example that the server machine
could not be reached or that the call timed out---and
\texttt{clnt\_perror} prints it.  A call that receives no reply times out;
over UDP the runtime retransmits the request while it waits.

\section{XDR data types and routines}
\label{sec:l13-xdr-types}

XDR provides the basic types familiar from C, such as integers,
floating-point numbers, enumerations, Booleans and
structures.\source{P4L1, XDR data types, XDR routines; RFC 4506
\cite{rfc4506}.}  It adds the following.
\begin{description}
  \item[\texttt{const}] a symbolic constant, translated into a
    \texttt{\#define};
  \item[\texttt{hyper}] a 64-bit integer;
  \item[\texttt{quadruple}] a 128-bit floating-point number;
  \item[\texttt{opaque}] uninterpreted binary data, represented in C as
    \texttt{char} bytes.
\end{description}
Arrays and strings are declared as in \cref{tab:l13-xdr-types}.  A
variable-length array, whose maximum length is given in angle brackets,
becomes a C structure with a length field and a pointer to the elements.
The size of this structure does not depend on the maximum length; the
elements reside in separately allocated memory.  A string becomes a
pointer to a NUL-terminated C string in memory, but it is transmitted as a
length followed by the characters (\cref{sec:l13-encoding}).

\begin{table}[htbp]
  \centering
  \caption{XDR arrays and strings and their C representation generated by
    \texttt{rpcgen}.}
  \label{tab:l13-xdr-types}
  \small
  \begin{tabular}{@{}>{\raggedright}p{27mm}>{\raggedright}p{31mm}%
      >{\raggedright\arraybackslash}p{52mm}@{}}
    \toprule
    Kind & XDR declaration & C representation \\
    \midrule
    Fixed-length array & \texttt{int data[8];} & \texttt{int data[8];} \\
    \addlinespace
    Variable-length array & \texttt{int data<8>;}
      & \texttt{struct \{ u\_int data\_len; int *data\_val; \} data;} \\
    \addlinespace
    String & \texttt{string name<32>;} & \texttt{char *name;} \\
    \addlinespace
    Fixed-length opaque & \texttt{opaque id[16];} & \texttt{char id[16];} \\
    \addlinespace
    Variable-length opaque & \texttt{opaque id<16>;}
      & \texttt{struct \{ u\_int id\_len; char *id\_val; \} id;} \\
    \bottomrule
  \end{tabular}
\end{table}

Pointers are declared with an asterisk, as in C.  Such \emph{optional
data} is encoded as a Boolean that states whether the pointer is non-null,
followed by the pointed-to data if it is, so that a linked list is
serialized element by element \cite{rfc4506}.

\subsection{XDR routines}

The file \texttt{sum\_xdr.c} contains one XDR routine for every data type
of the interface.  The routine for \texttt{sum\_in} as generated by
\texttt{rpcgen} is the following.

\begin{lstlisting}[language=C, numbers=none]
bool_t
xdr_sum_in(XDR *xdrs, sum_in *objp)
{
  if (!xdr_array(xdrs,
      (char **)&objp->values.values_val,
      (u_int *)&objp->values.values_len, MAX_VALUES,
      sizeof(int), (xdrproc_t)xdr_int))
    return FALSE;
  return TRUE;
}
\end{lstlisting}

The same routine serves for marshalling, unmarshalling and freeing, and
the field \texttt{x\_op} of the XDR stream \texttt{xdrs} selects the
direction.  When it encodes, the routine writes the length and the
elements into the buffer; when it decodes, it reads them and allocates the
array if \texttt{values\_val} is null; when it frees, it releases this
memory.  Memory allocated while unmarshalling must eventually be released.
On the client, \texttt{xdr\_free} applies a routine in the freeing direction
to a result that is no longer needed.  In servers generated with
\texttt{-M}, the server stub calls a procedure written by the developer,
\texttt{sum\_prog\_1\_freeresult}, after the result has been returned to
the client; it typically calls \texttt{xdr\_free}.

\section{XDR encoding}
\label{sec:l13-encoding}

An RPC message on the wire consists of three parts.\source{P4L1, encoding,
XDR encoding; RFC 5531 \cite{rfc5531}; RFC 4506 \cite{rfc4506}.}
\begin{description}
  \item[Transport header] the headers of the transport protocol, such as
    TCP or UDP, and of the protocols below it.
  \item[RPC header] the identification of the call: a request identifier
    (the \emph{xid}), with which the client matches replies to requests,
    the program, version and procedure numbers, and authentication data.
  \item[Data] the arguments or the result, encoded into a byte stream
    according to their data types.
\end{description}
XDR thus comprises both the IDL and the encoding, the binary representation
of the data on the wire.  Its encoding rules are few.
\begin{itemize}
  \item Every item is encoded in a multiple of \SI{4}{\byte}; shorter data
    are padded with zero bytes.
  \item Multi-byte values are transmitted in big-endian order.
  \item Integers are represented in two's complement.
  \item Floating-point numbers are represented in the IEEE~754 format.
\end{itemize}
Each machine converts between its native representation and XDR when it
marshals and unmarshals, so neither side needs to know the representation
of the other.

\begin{example}
  \label{ex:l13-encoding}
  The client of \cref{sec:l13-sun-binding} calls \texttt{SUM\_PROC} with the
  values 7, $-2$ and 300.  XDR encodes this variable-length array in the
  \SI{16}{\byte} of \cref{fig:l13-marshal}, the length followed by the three
  elements:
  \begin{center}
    \texttt{00 00 00 03 \quad 00 00 00 07 \quad ff ff ff fe \quad
      00 00 01 2c}
  \end{center}
  The element $-2$ appears in two's complement as \texttt{ff ff ff fe} and
  300 as \texttt{00 00 01 2c}, whatever the byte order of the client.  In
  front of these bytes the RPC header carries the request identifier and
  the program, version and procedure numbers of the call.
\end{example}

Strings show the difference between the representation in memory and on
the wire.

\begin{example}
  The string \texttt{"kernel"}, declared as \texttt{string name<32>}, takes
  \SI{7}{\byte} in the memory of client and server: six characters and the
  terminating NUL.  XDR encodes it in \SI{12}{\byte}: the length 6 in
  \SI{4}{\byte}, the six characters, and two bytes of padding to reach a
  multiple of four, without a terminator:
  \begin{center}
    \texttt{00 00 00 06 \quad 6b 65 72 6e \quad 65 6c 00 00}
  \end{center}
\end{example}

\begin{keypoint}
  XDR is both an IDL and an encoding.  Every item is encoded in multiples
  of \SI{4}{\byte}, in big-endian order, with two's complement integers and
  IEEE floating point; variable-length data carry their length.  A message
  adds an RPC header, which identifies the request, program, version and
  procedure, to the transport header.
\end{keypoint}

\section{Java RMI}
\label{sec:l13-rmi}

Java offers a counterpart to RPC that matches its object-oriented
semantics.\source{P4L1, Java RMI; Wollrath et al.\ \cite{wollrath1996}.}

\begin{definition}[Remote method invocation]
  The invocation of a method of an object that resides in another address
  space, with the syntax and semantics of a local method invocation, is
  called \term{remote method invocation}[リモートメソッド呼び出し] (RMI).  In
  Java RMI the address spaces are those of Java virtual machines (JVMs),
  on the same machine or on different machines.
\end{definition}
\index{RMI|see{remote method invocation}}%
\index{Java RMI|see{remote method invocation}}

The IDL of Java RMI is Java itself, a language-specific IDL.  A remote
interface is a Java interface that extends \texttt{java.rmi.Remote}, and
each of its methods declares the exception \texttt{RemoteException}, the
error notification through which partial failures are reported.
The following remote interface declares the \texttt{sum} service, and the
client code beneath it calls the service.

\begin{lstlisting}[language=Java, numbers=none]
public interface Summation extends java.rmi.Remote {
  int sum(int[] values) throws java.rmi.RemoteException;
}

/* client */
Registry reg = LocateRegistry.getRegistry("server");
Summation s = (Summation) reg.lookup("Summation");
int total = s.sum(new int[] { 7, -2, 300 });
\end{lstlisting}

The server creates an object that implements the interface and binds it
under a name in the \emph{RMI registry}; a client looks the name up and
obtains a stub for the remote object, whose methods it then invokes
(\cref{fig:l13-rmi}).  Arguments that are ordinary objects, such as the
array of integers above, are passed by value: they are serialized
together with the objects they refer to.  Remote objects are passed by
reference, as stubs.

\begin{figure}[htbp]
  \centering
  % Layers of Java RMI: stub and skeleton, remote reference layer, transport;
% the RMI registry maps names to remote object references.
% Usage: % Layers of Java RMI: stub and skeleton, remote reference layer, transport;
% the RMI registry maps names to remote object references.
% Usage: % Layers of Java RMI: stub and skeleton, remote reference layer, transport;
% the RMI registry maps names to remote object references.
% Usage: \input{figures/l13-rmi}   (width ~110 mm)
\begin{tikzpicture}[x=1mm, y=1mm, font=\footnotesize\sffamily,
  >={Stealth[length=1.6mm]},
  jvm/.style={draw=ln-rule, rounded corners=2pt, fill=ln-box},
  layer/.style={draw, fill=white, align=center, minimum width=36mm,
                minimum height=6.5mm, inner sep=1pt},
  stubl/.style={layer, fill=ln-accent!15},
  refl/.style={layer, fill=ln-accent!25},
  trl/.style={layer, fill=ln-accent!35},
  reg/.style={draw, rounded corners=2pt, fill=white, align=center,
              minimum width=26mm, minimum height=6.5mm, inner sep=1pt},
  hdr/.style={font=\scriptsize\sffamily\bfseries, text=ln-muted},
  lab/.style={font=\scriptsize\sffamily, text=ln-muted}]
  \draw[jvm] (0,0) rectangle (42,44);
  \draw[jvm] (68,0) rectangle (110,44);
  \node[hdr] at (21,2.6) {\iflangja{クライアント JVM}{client JVM}};
  \node[hdr] at (89,2.6) {\iflangja{サーバ JVM}{server JVM}};
  % client side
  \node[layer] (c0) at (21,37) {\iflangja{クライアント}{client object}};
  \node[stubl] (c1) at (21,28) {\iflangja{スタブ}{stub}};
  \node[refl]  (c2) at (21,19) {\iflangja{リモート参照層}{remote reference layer}};
  \node[trl]   (c3) at (21,10) {\iflangja{トランスポート層}{transport layer}};
  % server side
  \node[layer] (s0) at (89,37) {\iflangja{遠隔オブジェクト}{remote object}};
  \node[stubl] (s1) at (89,28) {\iflangja{スケルトン}{skeleton}};
  \node[refl]  (s2) at (89,19) {\iflangja{リモート参照層}{remote reference layer}};
  \node[trl]   (s3) at (89,10) {\iflangja{トランスポート層}{transport layer}};
  \draw[->] (c0) -- (c1);
  \draw[->] (c1) -- (c2);
  \draw[->] (c2) -- (c3);
  \draw[->] (s3) -- (s2);
  \draw[->] (s2) -- (s1);
  \draw[->] (s1) -- (s0);
  \draw[<->, thick, ln-accent] (c3) -- (s3)
    node[midway, below=1mm, lab] {\iflangja{ネットワーク}{network}};
  % registry
  \node[reg] (reg) at (55,52) {\iflangja{RMI レジストリ}{RMI registry}};
  \draw[->] ([xshift=10mm]s0.north) |- (reg.east)
    node[pos=0.25, right, lab] {\texttt{bind}};
  \draw[<->] ([xshift=-10mm]c0.north) |- (reg.west)
    node[pos=0.25, left, lab] {\texttt{lookup}};
\end{tikzpicture}
   (width ~110 mm)
\begin{tikzpicture}[x=1mm, y=1mm, font=\footnotesize\sffamily,
  >={Stealth[length=1.6mm]},
  jvm/.style={draw=ln-rule, rounded corners=2pt, fill=ln-box},
  layer/.style={draw, fill=white, align=center, minimum width=36mm,
                minimum height=6.5mm, inner sep=1pt},
  stubl/.style={layer, fill=ln-accent!15},
  refl/.style={layer, fill=ln-accent!25},
  trl/.style={layer, fill=ln-accent!35},
  reg/.style={draw, rounded corners=2pt, fill=white, align=center,
              minimum width=26mm, minimum height=6.5mm, inner sep=1pt},
  hdr/.style={font=\scriptsize\sffamily\bfseries, text=ln-muted},
  lab/.style={font=\scriptsize\sffamily, text=ln-muted}]
  \draw[jvm] (0,0) rectangle (42,44);
  \draw[jvm] (68,0) rectangle (110,44);
  \node[hdr] at (21,2.6) {\iflangja{クライアント JVM}{client JVM}};
  \node[hdr] at (89,2.6) {\iflangja{サーバ JVM}{server JVM}};
  % client side
  \node[layer] (c0) at (21,37) {\iflangja{クライアント}{client object}};
  \node[stubl] (c1) at (21,28) {\iflangja{スタブ}{stub}};
  \node[refl]  (c2) at (21,19) {\iflangja{リモート参照層}{remote reference layer}};
  \node[trl]   (c3) at (21,10) {\iflangja{トランスポート層}{transport layer}};
  % server side
  \node[layer] (s0) at (89,37) {\iflangja{遠隔オブジェクト}{remote object}};
  \node[stubl] (s1) at (89,28) {\iflangja{スケルトン}{skeleton}};
  \node[refl]  (s2) at (89,19) {\iflangja{リモート参照層}{remote reference layer}};
  \node[trl]   (s3) at (89,10) {\iflangja{トランスポート層}{transport layer}};
  \draw[->] (c0) -- (c1);
  \draw[->] (c1) -- (c2);
  \draw[->] (c2) -- (c3);
  \draw[->] (s3) -- (s2);
  \draw[->] (s2) -- (s1);
  \draw[->] (s1) -- (s0);
  \draw[<->, thick, ln-accent] (c3) -- (s3)
    node[midway, below=1mm, lab] {\iflangja{ネットワーク}{network}};
  % registry
  \node[reg] (reg) at (55,52) {\iflangja{RMI レジストリ}{RMI registry}};
  \draw[->] ([xshift=10mm]s0.north) |- (reg.east)
    node[pos=0.25, right, lab] {\texttt{bind}};
  \draw[<->] ([xshift=-10mm]c0.north) |- (reg.west)
    node[pos=0.25, left, lab] {\texttt{lookup}};
\end{tikzpicture}
   (width ~110 mm)
\begin{tikzpicture}[x=1mm, y=1mm, font=\footnotesize\sffamily,
  >={Stealth[length=1.6mm]},
  jvm/.style={draw=ln-rule, rounded corners=2pt, fill=ln-box},
  layer/.style={draw, fill=white, align=center, minimum width=36mm,
                minimum height=6.5mm, inner sep=1pt},
  stubl/.style={layer, fill=ln-accent!15},
  refl/.style={layer, fill=ln-accent!25},
  trl/.style={layer, fill=ln-accent!35},
  reg/.style={draw, rounded corners=2pt, fill=white, align=center,
              minimum width=26mm, minimum height=6.5mm, inner sep=1pt},
  hdr/.style={font=\scriptsize\sffamily\bfseries, text=ln-muted},
  lab/.style={font=\scriptsize\sffamily, text=ln-muted}]
  \draw[jvm] (0,0) rectangle (42,44);
  \draw[jvm] (68,0) rectangle (110,44);
  \node[hdr] at (21,2.6) {\iflangja{クライアント JVM}{client JVM}};
  \node[hdr] at (89,2.6) {\iflangja{サーバ JVM}{server JVM}};
  % client side
  \node[layer] (c0) at (21,37) {\iflangja{クライアント}{client object}};
  \node[stubl] (c1) at (21,28) {\iflangja{スタブ}{stub}};
  \node[refl]  (c2) at (21,19) {\iflangja{リモート参照層}{remote reference layer}};
  \node[trl]   (c3) at (21,10) {\iflangja{トランスポート層}{transport layer}};
  % server side
  \node[layer] (s0) at (89,37) {\iflangja{遠隔オブジェクト}{remote object}};
  \node[stubl] (s1) at (89,28) {\iflangja{スケルトン}{skeleton}};
  \node[refl]  (s2) at (89,19) {\iflangja{リモート参照層}{remote reference layer}};
  \node[trl]   (s3) at (89,10) {\iflangja{トランスポート層}{transport layer}};
  \draw[->] (c0) -- (c1);
  \draw[->] (c1) -- (c2);
  \draw[->] (c2) -- (c3);
  \draw[->] (s3) -- (s2);
  \draw[->] (s2) -- (s1);
  \draw[->] (s1) -- (s0);
  \draw[<->, thick, ln-accent] (c3) -- (s3)
    node[midway, below=1mm, lab] {\iflangja{ネットワーク}{network}};
  % registry
  \node[reg] (reg) at (55,52) {\iflangja{RMI レジストリ}{RMI registry}};
  \draw[->] ([xshift=10mm]s0.north) |- (reg.east)
    node[pos=0.25, right, lab] {\texttt{bind}};
  \draw[<->] ([xshift=-10mm]c0.north) |- (reg.west)
    node[pos=0.25, left, lab] {\texttt{lookup}};
\end{tikzpicture}

  \caption{Java RMI.  Invocations pass from the client through the stub,
    the remote reference layer and the transport layer to the server, where
    the skeleton delivers them to the remote object.}
  \label{fig:l13-rmi}
\end{figure}

\begin{table}[htbp]
  \centering
  \caption{Design choices of Sun RPC and Java RMI.}
  \label{tab:l13-design}
  \small
  \begin{tabular}{@{}>{\raggedright}p{22mm}>{\raggedright}p{45mm}%
      >{\raggedright\arraybackslash}p{45mm}@{}}
    \toprule
    Choice & Sun RPC & Java RMI \\
    \midrule
    Binding & per-machine registry daemon, \texttt{rpcbind}
      & RMI registry maps names to remote objects \\
    \addlinespace
    IDL & XDR, language-agnostic, compiled by \texttt{rpcgen}
      & Java remote interfaces, language-specific \\
    \addlinespace
    Encoding & XDR encoding & Java object serialization \\
    \addlinespace
    Pointers and references & allowed; pointed-to data serialized
      & objects copied; remote objects passed as stubs \\
    \addlinespace
    Partial failures & timeouts and retries; error status in the client
      handle & \texttt{RemoteException} \\
    \bottomrule
  \end{tabular}
\end{table}

Java RMI has three layers on each side.  The topmost is the stub on the
client and, on the server, the \emph{skeleton}, which unmarshals an
invocation and dispatches it to the remote object.\margin{Since JDK~1.2 the
stub dispatches directly, and no skeleton is generated.}  The RMI runtime
provides the two layers beneath them.
\begin{description}
  \item[Remote reference layer] It determines the invocation semantics:
    \emph{unicast}, to a single object, or \emph{broadcast}, to a group of
    objects such as replicas, returning the first response or returning
    only if all responses match.
  \item[Transport layer] It sets up and manages the connections and
    transfers the data.  The architecture allows different transports, such
    as TCP, UDP or shared memory.
\end{description}
\Cref{tab:l13-design} compares the design choices of Sun RPC and Java RMI.

\needspace{13\baselineskip}
\begin{keypoint}[Lesson summary]
  RPC lets a process call a procedure in another address space, usually on
  another machine, with the synchronous semantics of a local call.  Stubs
  generated from an IDL marshal and unmarshal arguments and results, and
  the RPC runtime transfers the messages.  Every RPC system must decide how
  clients bind to servers through a registry, how interfaces are described
  and data encoded, how pointer arguments are treated, and how partial
  failures, which a client cannot diagnose, are reported.  Sun RPC uses
  XDR as IDL and encoding, \texttt{rpcgen} to generate stubs, and a
  per-machine \texttt{rpcbind} registry; Java RMI uses Java as IDL and adds
  a remote reference layer that defines invocation semantics.
\end{keypoint}

\end{document}
